% !TEX root = ../main.tex
\chapter{구현과 실증 결과}
\label{ch:results}

이 장에서는 \ref{ch:methodology}장의 설계를 아비트럼 원(Arbitrum One)에서 하나의 사례 연구로 실제로 돌려 봐요. 체인은 하나이고, 점수 기간(점수를 매길 자료를 모은 기간)은 동결(점수를 고정하는 일)마다 하나예요. 같은 기간 비교에는 모두 스마트 계약(블록체인 위에서 저절로 실행되는 프로그램)의 매칭 코호트(모든 점수가 똑같이 점수를 매기는 계약 묶음) 하나를 써요. 시간 홀드아웃(나중 기간을 떼어 두고 맞혀 보는 시험) W0과 W1에는 스펜더(spender, 허락을 받아 대신 쓰는 쪽) 코호트와 트레이더(거래하는 쪽) 코호트를 써요. 크기에 따른 실행 시간을 볼 때는 매칭 코호트에서 일부만 뽑은 부분 표본을 써요. 이 장은 먼저 한 층만 걷는 두 걷기를 보고해요. 두 걷기의 비용, 기간 안에서의 일치, 순위 갈라짐, 강건성 점검(조건을 바꿔도 결과가 버티는지 보기), 홀드아웃을 차례로 보고해요(\ref{sec:benchmark-results}절부터 \ref{sec:holdout-results}절까지). 그다음에는 이 논문의 방법인 결합 걷기 C-PR(결합 페이지랭크, 두 그래프를 함께 걷는 점수)을, 그 규칙 아래에서, 그리고 다른 모든 점수와 나란히 보고해요(\ref{sec:cpr-results}절). 이어서 두 종류의 화살표(엣지) 어느 쪽으로도 만들지 않은 결과 값에서의 비교가 나와요(\ref{sub:neutral-label-results}절). 이 장은 \ref{sec:amount-revision}절의 수정 분석(등록 뒤에 금액을 달러로 바꿔 다시 한 분석)을 보고해요. 수정 분석에서는 금액을 달러(USD)로 매기고, 다섯 토큰을 쓰고, 각 동결 전 열두 달의 자료로 점수를 매겨요. 이 분석은 등록 분석(미리 등록한 그대로 한 분석)의 결과를 안 뒤에 정의했어요. 그래서 이 장의 어떤 결과도 등록 분석의 결과가 아니에요. 규칙~A와~B, R1--R3 모두를 기준값과 문구까지 그대로, 고친 점수에 재분석으로 적용해요. 등록 분석도 같은 판정에 이르렀지만, R1만은 달랐어요. 등록 분석에서는 R1이 충족되지 않았어요(부록~\ref{sec:map-registered}, 표~\ref{tab:registered-vs-revised}). 모든 표는 \ref{sec:reproducibility}절의 평가 요약 파일에서 만들어요. 켄달의 타우 $\tau$(두 줄 세우기가 얼마나 같은 순서인지 나타내는 수)에는 모두 신뢰구간(믿을 만한 범위)이 함께 나와요. 다만 금액 기울기 점검과 표본 정의 점검에는 구간이 없어요. 그리고 페이지랭크(PageRank)는 모두 그것이 매끄럽게 다듬는(평활하는) 원시 차수(화살표를 그냥 센 수)와 견주어 읽어요.

\dissertationSubheading[sec:dataset]{사례 연구 배경: 아비트럼 원}

아비트럼 원을 고른 까닭은 세 가지예요. 첫째, 아비트럼 원은 많이 쓰이는 이더리움 롤업(레이어 2)이고, 그 위에서 충분히 자리 잡은 디파이(DeFi, 탈중앙 금융) 서비스(프로토콜)들이 돌아가요(\ref{sec:ethereum-l2}절). 둘째, 아비트럼 원의 로그(계약이 남기는 기록)는 공개된 빅쿼리(BigQuery) 내보내기 자료에 들어 있고, 처리 과정(파이프라인)이 바로 이 자료를 읽어요. 셋째, 아비트럼 원의 증거금 거래(빌린 돈을 보태 하는 거래) 서비스 가운데 하나인 GMX~V2가 포지션 닫기를 기록해요. 트레이더 라벨(나중에 채점에 쓰는 정답 값)은 이 기록에서 나와요. 각 점수 기간은 그 동결 전 열두 달을 덮어요. 같은 기간 평가와 W1에서는 2025년 7월 1일부터 2026년 6월 30일까지이고, W0에서는 2025년 4월 1일부터 2026년 3월 31일까지예요. 로지스틱 감쇠(오래된 일일수록 가볍게 세는 무게)는 동결 시점을 기준으로 해요(그림~\ref{fig:logistic-decay-methods}). 그래서 AWP(적응형 가중 페이지랭크)와 엔도스랭크(EndorseRank)에서 화살표 무게의 대부분은 기간의 가장 최근 여섯 달이 차지해요.

\dissertationSubsubheading[sub:cohorts]{코호트}

계약 집합 네 개를 써요. 이 집합에 든 주소는 모두 계약 코드를 가지고 있어요. 외부 소유 계정(EOA, 개인 키로 움직이는 계정)은 그래프(점과 화살표로 이은 그림)에 주인(허락을 해 주는 쪽), 보낸 쪽(송신자), 받는 쪽(수신자)으로만 들어가고, 점수를 받지 않아요(\ref{sec:data-sources}절).

\begin{itemize}
\item 매칭 코호트($n=6{,}483$). $T_{\mathrm{obs}}$의 점수 기간에, 다섯 토큰 가운데 하나의 0이 아닌 \texttt{Approval}을 서로 다른 주인 세 명 이상에게서 받은 계약들이에요(\ref{sec:data-sources}절). 이 계약들은 분석 계획의 코호트인 계약 27,844개의 일부예요. 계획의 코호트는 2023년 10월 1일부터 모든 토큰에 같은 규칙을 적용했어요. 그리고 후보 29,022개 가운데 코드가 없는 1,121개와 EIP-7702 위임 표시를 가진 57개를 EOA로 보고 뺐어요. 코호트에 드는지는 어느 그래프에도 달려 있지 않아요. 코호트 계약 가운데 $191$개에는 허락(allowance, 토큰을 얼마까지 대신 써도 된다는 허락) 화살표가 하나도 닿지 않아요($T_{\mathrm{obs}}$에는, 기간 안에 그 계약들이 해 주었거나 그 계약들이 받은 허락이 모두 0이었어요). $1{,}192$개에는 송금(transfer) 화살표가 하나도 닿지 않고, $30$개에는 어느 쪽 화살표도 닿지 않아요. 어떤 그래프에 없는 계약은 그 방법에서 0점을 받아요. 같은 기간 일치도, 차이 검정, 벤치마크(속도 재기)는 모두 이 코호트에서 계산해요. W0의 자료는 W0만의 코호트를 둘러싸고 뽑았어요. 이 코호트는 W0의 점수 기간에서 같은 규칙을 만족하는 계약 $6{,}343$개이고, 그 가운데 $5{,}383$개는 매칭 코호트에도 들어 있어요.
\item W0과 W1의 스펜더 코호트($n=7{,}914$와 $7{,}960$). W0에서는 점수를 $t_1$(2026년 3월 31일, 23:59:59~UTC)에 동결해요. 이 코호트는 $t_1$에 다섯 토큰 가운데 하나의 양수인 마지막 허락(가장 최근 허락)을 하나 이상 가지고 있고 계약 코드를 가진 모든 스펜더예요. 그 허락은 W0의 점수 기간 안에 정해진 것이어야 해요. 이런 허락을 가진 스펜더 7,962개 가운데 7,914개예요. W1에서는 점수를 $T_{\mathrm{obs}}$에 동결하고, 같은 규칙으로 스펜더 계약 7,960개가 나와요. 스펜더 코호트는 매칭 코호트의 부분 집합이 아니에요. 스펜더는 양수 허락을 해 준 주인이 한 명만 있으면 되기 때문이에요. 이 계약들의 점수는, 각 동결 뒤 석 달 안에 처음 나타난 주인--스펜더 허락 쌍과 송금을 보낸 쪽에 견주어 시험해요(\ref{sec:holdout-protocol}절).
\item 트레이더 코호트($n=87$과 $70$). 계약 코드를 가지고 있고, 동결일 그래프의 점(노드)이고, 라벨 기간 안에 포지션을 세 개 이상 닫은 GMX~V2 계정들이에요. W0에서는 87개이고, 그 가운데 39개는 점수 기간에 허락을 받은 적이 있었어요. W1에서는 70개이고, 그 가운데 37개가 그랬어요. 계약 계정들은 2023년 10월부터 2026년 6월까지 GMX~V2 포지션 133,685개를 닫았어요. 하지만 그 가운데 마지막 여섯 달에는 9,809개만 닫았고, 2026년 7월부터 9월까지는 6,805개만 닫았어요(부록~\ref{sec:data-gmx}의 표~\ref{tab:gmx-quarters}). 그래서 트레이더 결과는 모두 적은 수의 계약에 기대고, 구간이 넓어요.
\item 매칭 코호트의 부분 표본($n=1{,}000$, $2{,}000$, $5{,}000$). SHA256으로 늘 똑같이 뽑히도록 정한 부분 표본(시드 \texttt{contract-benchmark-v1})이고, 그다음은 코호트 전체예요. 이 표본들은 크기에 따른 실행 시간과 표본 정의에 대한 민감도(조건을 바꿔 결과가 얼마나 움직이는지 보기)를 보는 데에만 써요(\ref{sub:expanded-pool}절). 각 단계는 뽑힌 계약에 닿는, 뽑아 낸 화살표를 남겨요. 그래서 단계의 화살표 수는 $n$과 함께 늘어나요.
\end{itemize}

$T_{\mathrm{obs}}$에서 허락 그래프에는 점 2,514,846개와 화살표 3,808,949개가 있어요. 송금 그래프에는 점 3,603,300개와 화살표 8,426,124개가 있어요. 표~\ref{tab:benchmark-runtime}의 풀이 프로그램(solver)이 쓰는 그래프는 각각 주소 864개와 1,573개를 빼요. 이 주소들이 낀 쌍은 모두 부동소수점 계산(컴퓨터가 소수를 다루는 방식)에서 금액 무게가 0이에요. 같은 계약들 위에서 허락 그래프의 화살표 수는 송금 그래프의 45.2\%예요. 주인--스펜더 쌍은 다섯 토큰 가운데 하나의 마지막 허락이 양수인 동안에만 허락 화살표가 되기 때문이에요. 그 허락은 기간 안에 정해진 것이어야 해요. 이와 달리 기간 안에 이 토큰들의 0이 아닌 송금을 주고받은 주소 쌍은 모두 송금 화살표가 돼요. 표본을 뽑는 방식은 두 그래프에서 같아요. 그래서 점수들 사이의 차이는 화살표와 다시 출발(재시작)에서 비롯돼요.

뽑아 낸 행(기록 한 줄)의 대부분은 다섯 토큰의 것이에요. $T_{\mathrm{obs}}$의 점수 기간에 있는, 계획의 코호트의 에고 네트워크(코호트 계약에 닿는 화살표만 모은 그래프) 이벤트 가운데, 수정 분석은 허락 85,986,562개(63.3\%)와 송금 442,722,039개(76.2\%)를 남겨요. 그리고 남긴 송금에는 모두 가격이 있어요. 부록~\ref{sec:data-tokens}의 표~\ref{tab:usd-tokens}에는 이 수치가 토큰마다 있어요.

\dissertationSubsubheading[sub:implementation]{구현}

모든 알고리즘은 페이지랭크를 하나의 파이썬(Python) 구현으로 계산해요. 감쇠 계수(화살표를 따라갈 확률)는 $d=0.85$이고, 허용 오차는 $\varepsilon=10^{-8}$이고, 반복(한 바퀴) 상한은 $I_{\max}=300$이에요. 엔도스랭크와 AWP는 화살표 목록과 재시작 벡터가 달라요. C-PR은 두 화살표 목록을 점마다 섞어요(식~\eqref{eq:coupled-transition}). 평가 실행은 한 번마다 점수, 대리 지표(대신 재는 값)별 통계, 부트스트랩(다시 뽑기) 구간을 요약 파일 하나에 써요. 벤치마크는 잰 시간을 다른 파일에 써요. 이 장의 표는 2026년 10월 28일 09:16~UTC부터 돌린 실행에서 나왔어요. $d=0.85$에서 모든 풀이는 반복 상한보다 훨씬 일찍 수렴했어요(값이 더는 거의 바뀌지 않게 되었어요). 부록~\ref{ch:appendix-eval}에는 실행 목록, 저장소 버전, 설정 파일, 요약 파일이 실려 있어요.

\dissertationSubheading[sec:benchmark-results]{계산 벤치마크(주장~C5)}

주장~C5(이 연구가 내세우는 것)는 같은 입력에서 각 점수를 계산하는 데 얼마나 드는지에 대한 것이에요. 표~\ref{tab:benchmark-runtime}은 \ref{sec:computational-benchmark}절의 시간 재는 절차를 따라, 알고리즘마다 몇 가지를 보고해요. 대상은 매칭 코호트의 전체 그래프예요. 보고하는 것은 다음과 같아요. 시간을 재지 않는 준비 실행을 한 번 한 뒤 시간 잰 실행을 다섯 번 해서 얻은 평균 푸는 시간과 그 표준편차(SD, 값이 퍼진 정도), 준비 실행 풀이에서 가장 많이 쓴 메모리, 수렴할 때까지의 반복 횟수, 그리고 풀이 프로그램이 쓰는 그래프의 점과 화살표 수예요. 각 그래프와 설정은 한 번씩만 시간을 쟀어요. 그래서 같은 측정이 표~\ref{tab:benchmark-scaling}의 코호트 전체 행과 표~\ref{tab:robustness-damping}의 $d=0.85$ 행에 다시 나와요. 부록~\ref{sec:benchmark-run-logs}의 표~\ref{tab:benchmark-runs}에는 시간 잰 실행이 모두 실려 있어요.

% Generated by experiments/scripts/export_latex.py -- do not edit by hand.
% Source: experiments/data/2-processed-tables-and-evaluations/revised-usd-12m/benchmark/benchmark.json
\begin{table}[htbp]
\centering
\caption[매칭 코호트의 전체 그래프에서 잰 계산 벤치마크.]{매칭 코호트의 전체 그래프에서 잰 계산 벤치마크(속도 재기)예요(계약 $n=6{,}483$개. AWP 대 엔도스랭크의 실행 시간 비율은 3.2$\times$이고, 이때 화살표 수의 비율은 2.2$\times$예요). 실행 시간은 시간을 재지 않는 준비 실행을 한 번 한 뒤, 시간을 잰 풀이 프로그램(solver) 호출 5번에서 실제 걸린 시간(초)의 평균이에요. 최대 메모리는 그 준비 실행에서만 재요. 표준편차(SD)는 같은 5번의 실행에서 구했어요. 시간을 잰 호출 한 번은 화살표 목록을 점수로 바꿔요. 여기에는 점에 번호 붙이기, 듬성듬성한 행렬 만들기, 행 정규화(각 행의 합을 1로 맞추기), 재시작 벡터 만들기, 거듭제곱 반복이 들어가요. $|V|$와 $|E|$는 각 방법의 풀이 프로그램이 쓰는 그래프의 점 수와 서로 다른 화살표 수예요. C-PR은 허락 층과 송금 층을 합친 그래프 위를 걸어요. 그래서 C-PR의 $|E|$는 두 층의 서로 다른 쌍을 모두 세요. S-PR은 송금 층에서 돌아가요. S-PR의 실행 시간과 반복 수에는, 재시작 벡터를 만들어 주는 허락 층을 푸는 계산도 들어 있어요. 모든 실행은 표~\ref{tab:benchmark-runs}에 있어요.}
\label{tab:benchmark-runtime}
\fitwidth{%
\begin{tabular}{lrrrrrr}
\toprule
방법 & 실행 시간 (초) & 표준편차 (초) & 최대 메모리 (MB) & 반복 수 & $|V|$ & $|E|$ \\
\midrule
EndorseRank & 2.657 & 0.109 & 325.3 & 57 & 2{,}513{,}982 & 3{,}808{,}949 \\
AWP & 8.530 & 0.618 & 663.7 & 94 & 3{,}601{,}727 & 8{,}426{,}124 \\
\midrule
C-PR ($\lambda=0.5$) & 11.563 & 0.214 & 719.1 & 82 & 4{,}111{,}490 & 10{,}562{,}390 \\
S-PR & 10.883 & 0.330 & 509.4 & 151 & 3{,}603{,}300 & 8{,}432{,}635 \\
\bottomrule
\end{tabular}
}
\end{table}


엔도스랭크는 허락 그래프를 $2.66$\,초에 풀고, AWP는 송금 그래프를 $8.53$\,초에 풀어요. 그래서 AWP는 $3.21$배 오래 걸리고, 이때 화살표 수 비율은 $2.21$이에요. 거듭제곱 반복(같은 계산을 되풀이하기)에는 한 번마다, 화살표 $|E|$개에 걸친 듬성듬성한(희소) 행렬 곱셈 한 번과 점 $|V|$개에 걸친 벡터 계산 몇 번이 들어요(부록~\ref{sec:pr-sparsity}). 행렬을 조립하는 시간까지 넣어 반복 횟수로 나누면, AWP의 풀이는 반복 한 번에 $0.091$\,초, 엔도스랭크의 풀이는 $0.047$\,초가 걸려요. AWP가 $1.95$배 오래 걸리는 것이고, 이 값은 점 수 비율 $1.43$과 화살표 수 비율 사이에 있어요. 그리고 AWP는 반복이 $1.65$배 많이 필요해요($94$번 대 $57$번). 이 두 인수(곱해지는 값)를 합치면 시간 비율이 나와요. 같은 그래프에서 이 비율은 $d=0.75$, $0.85$, $0.95$일 때 $3.00$, $3.21$, $2.56$이에요(\ref{sub:damping}절). 그래서 이 비율은 화살표 수 비율을 따르지 않아요. 두 점수는 같은 풀이 프로그램을 써요. 그래서 속도 차이는 각 그래프의 크기와, 그 그래프에서 걷기에 필요한 반복 횟수에서 나와요. 벤치마크는 이 둘을 갈라 보지 않아요. 이 차이는 엔도스랭크 알고리즘의 장점이 아니에요.

C-PR은 두 화살표 집합을 합친 것(서로 다른 쌍 $10{,}562{,}390$개, 점 $4{,}111{,}490$개) 위를 걸어요. 하지만 C-PR의 풀이 프로그램은 층마다 행렬을 하나씩 두고, 반복할 때마다 두 행렬을 모두 곱해요. 저장된 칸은 모두 합쳐 약 $12.2$백만 개에 이르고, 이것은 AWP 화살표 수의 $1.45$배예요. C-PR은 더 적은 반복으로 수렴하고($82$번 대 $94$번), $11.56$\,초가 걸려요. AWP보다 $1.36$배 오래 걸리는 거예요. S-PR(보증 씨앗 페이지랭크)은 송금 층 위에서 돌고, $10.88$\,초가 걸려요. S-PR의 시간과 반복 횟수에는 재시작 벡터를 대 주는 허락 층의 풀이도 들어 있어요. 반복 $151$번은 그 풀이의 $56$번과, 씨앗을 쓴 송금 걷기의 $95$번을 더한 거예요. 송금 층을 걷는 알고리즘은, 혼자 걷든 허락 화살표와 함께 걷든, 모두 그 층의 크기만큼 값을 치러요.

이 결과에는 주의할 점이 두 가지 있어요. 이 시간은 미리 만들어 둔 화살표 목록에서 풀이 프로그램을 부르는 시간만 담아요. 처음부터 모두 새로 계산하려면 읽어 들이기(파싱), 해독(디코딩), 화살표 목록 만들기가 더해져요. 이 일들은 모든 알고리즘이 전부 또는 일부를 함께 해요. 그리고 더 빠른 방법이 꼭 더 좋은 방법은 아니에요.

\dissertationSubsubheading[sub:runtime-scaling]{매칭 코호트 부분 표본에서의 실행 시간}

표~\ref{tab:benchmark-scaling}에는 \ref{sub:cohorts}절의 각 단계에서 잰 엔도스랭크와 AWP의 실행 시간, 그 비율, 화살표 수 비율, 그리고 두 풀이 프로그램 그래프의 점과 화살표 수가 있어요.

% Generated by experiments/scripts/export_latex.py -- do not edit by hand.
% Source: experiments/data/2-processed-tables-and-evaluations/revised-usd-12m/benchmark/benchmark.json
\begin{table}[htbp]
\centering
\caption[평가 집합의 크기에 따른 실행 시간.]{평가 집합의 크기에 따라 엔도스랭크(ER)와 AWP의 실행 시간이 어떻게 바뀌는지 보여 줘요. 평가 집합은 매칭 코호트에서 SHA256으로 늘 같은 결과가 나오게 뽑은 부분 표본이에요(시드 \texttt{contract-benchmark-v1}). 가장 큰 것은 계약 $n=6{,}483$개인 코호트 전체예요. 단계마다 뽑힌 계약에 닿는, 추출한 화살표만 남겨요. 그래서 두 풀이 프로그램 그래프의 $|V|$와 $|E|$를 단계마다 적었어요. AWP/ER 시간과 AWP/ER $|E|$는 같은 행에 있는 두 실행 시간의 비율과 두 화살표 수의 비율이에요. 실행 시간은 시간을 재지 않는 준비 실행을 한 번 한 뒤, 시간 잰 풀이 프로그램 호출 5번에서 실제 걸린 시간을 초 단위로 평균한 값이에요. $^{a}$코호트 전체 행은 표~\ref{tab:benchmark-runtime}의 측정값을 다시 쓴 것이고, 시간을 다시 재지 않았어요.}
\label{tab:benchmark-scaling}
\fitwidth{%
\begin{tabular}{rrrrrrrrr}
\toprule
$n$ & ER (초) & AWP (초) & AWP/ER 시간 & AWP/ER $|E|$ & ER $|V|$ & ER $|E|$ & AWP $|V|$ & AWP $|E|$ \\
\midrule
1{,}000 & 0.498 & 1.307 & 2.6$\times$ & 2.8$\times$ & 451{,}248 & 485{,}077 & 953{,}795 & 1{,}382{,}177 \\
2{,}000 & 0.318 & 2.383 & 7.5$\times$ & 2.9$\times$ & 760{,}410 & 897{,}812 & 1{,}459{,}514 & 2{,}647{,}854 \\
5{,}000 & 1.672 & 6.620 & 4.0$\times$ & 2.7$\times$ & 1{,}834{,}957 & 2{,}577{,}926 & 3{,}183{,}096 & 7{,}035{,}600 \\
6{,}483$^{a}$ & 2.657 & 8.530 & 3.2$\times$ & 2.2$\times$ & 2{,}513{,}982 & 3{,}808{,}949 & 3{,}601{,}727 & 8{,}426{,}124 \\
\bottomrule
\end{tabular}
}
\end{table}


\ref{sub:expanded-pool}절의 단계들은 크기뿐 아니라 화살표 수 비율도 바꿔요. 송금 그래프의 화살표 수는 허락 그래프의 $2.85$배($n=1{,}000$), $2.95$배($n=2{,}000$), $2.73$배($n=5{,}000$)이고, 코호트 전체에서는 $2.21$배예요. 풀이 시간이 $|E|$만 따라간다면, 시간 비율도 이 화살표 수 비율을 따라갈 거예요. 실제 시간 비율은 $2.63$, $7.49$, $3.96$, $3.21$이에요. 어느 단계에서도 시간 비율은 화살표 수 비율 가까이에 머물지 않아요. $n=1{,}000$에서는 화살표 수 비율의 $0.92$배예요. 이때 AWP는 반복이 $96$번, 엔도스랭크는 $43$번 필요해요. $n=2{,}000$에서는 $2.54$배이고, AWP는 $94$번, 엔도스랭크는 $30$번이 필요해요. $n=5{,}000$과 코호트 전체에서는 $1.45$배예요. 이때 엔도스랭크는 $57$번, AWP는 $95$번과 $94$번이 필요해요. 엔도스랭크의 풀이는 $0.32$\,초에서 $2.66$\,초가 걸려요. 비율을 정하는 것은 화살표 수만이 아니라, 각 부분 그래프에서 각 걷기에 필요한 반복 횟수예요. 단계들은 코호트 그래프 하나에서 뽑은 부분 표본이에요. 그래서 단계들은 전체 그래프에 이르기까지 화살표에 따라 실행 시간이 어떻게 늘어나는지 보여 주지만, 그보다 큰 그래프에 대해서는 보여 주지 않아요. 그림~\ref{fig:runtime-scaling}에는 단계의 크기에 따라 두 실행 시간을 그렸고, 점 옆에 각 그래프의 화살표 수를 적었어요.

\begin{figure}[htbp]
\centering
\includegraphics[width=0.9\textwidth]{\figpath{runtime-scaling.pdf}}
\caption[{뽑은 계약 수에 따른 엔도스랭크와 AWP의 실행 시간.}]{표~\ref{tab:benchmark-scaling}의 각 단계에 든 계약 수에 따라, 엔도스랭크와 AWP의 풀이 프로그램을 한 번 부를 때 실제 걸린 시간을 그렸어요. 두 축은 모두 로그 눈금이에요. 점 옆의 글자는 그 그래프의 화살표 수 $|E|$예요. 오차 막대는 시간 잰 실행 다섯 번의 표준편차예요.}
\label{fig:runtime-scaling}
\end{figure}

\dissertationSubheading[sec:same-window-alignment]{같은 기간 일치도(주장~C4)}

주장~C4는 각 점수가, 그 점수를 계산한 기간 안에서 무엇과 잘 맞는지 물어요. 표~\ref{tab:alignment-hybrid}은 모든 점수의 지표 가족(비슷한 지표 묶음) 평균 켄달의 $\tau$를 송금, 허락, 시빌(가짜 주소를 많이 만들어 속이는 공격) 안정성 가족마다 보고해요. 값마다 95\% 짝지은 부트스트랩 신뢰구간이 함께 나와요. 원시 차수에는 행이 없어요. 이 기간에서는 원시 차수가 대리 지표이기 때문이에요. 송금 입차수(송금을 보내온 서로 다른 주소 수)는 송금 가족에 들고, 받은 허락 수(허락 입차수)는 허락 가족에 들어요. 규칙~A의 보조 부분은 분석 계획에서 정해 두었고(\ref{sec:holdout-protocol}절), 여기서는 고친 점수에 재분석으로 적용해요. 이것을 빼면, 아래의 비교 가운데 어느 것도 미리 정해 두지 않았어요(\ref{ch:methodology}장).

% Generated by experiments/scripts/export_latex.py -- do not edit by hand.
% Source: experiments/data/2-processed-tables-and-evaluations/revised-usd-12m/same-window/eval_summary.json
\begin{table}[htbp]
\centering
\caption[이 연구의 모든 점수에 대한 같은 기간 지표 가족 평균 켄달의 $\tau$.]{이 연구의 모든 점수에 대해, 같은 기간의 지표 가족 평균 켄달의 $\tau$와 부트스트랩 구간을 보여 줘요(매칭 코호트, 계약 $n=6{,}483$개). 점수는 엔도스랭크와 AWP, 다섯 가지 $\lambda$ 값의 결합 연산자 C-PR, 그리고 보증 씨앗 걷기 S-PR이에요. $\lambda=1$은 허락 층만 걷는 것이에요(달러로 나타낸 마지막 허락, 각각 상한을 둠). $\lambda=0$은 송금 층만 걷는 것이에요(달러 금액을 AWP의 로지스틱 감쇠로 줄임). 둘 다 고르게 다시 출발해요. 원시 차수는 송금 가족과 허락 가족의 대리 지표이기 때문에 따로 행이 없어요. 활동 재시작을 쓰는 엔도스랭크는 표~\ref{tab:endorserank-restarts}에 있어요. 신뢰구간(CI)은 계약을 다시 뽑아 구한 95\% 백분위 부트스트랩 구간이에요(짝지은 다시 뽑기 400번, 시드 42).}
\label{tab:alignment-hybrid}
\fitwidth{%
\begin{tabular}{lccc}
\toprule
방법 & 송금 & 허락 & 시빌 \\
\midrule
EndorseRank & 0.196 [0.180, 0.216] & 0.560 [0.548, 0.571] & 0.026 [0.012, 0.041] \\
AWP & 0.736 [0.728, 0.744] & 0.226 [0.210, 0.245] & 0.419 [0.408, 0.434] \\
\midrule
C-PR ($\lambda=1$) & 0.196 [0.180, 0.216] & 0.561 [0.550, 0.572] & 0.028 [0.014, 0.042] \\
C-PR ($\lambda=0$) & 0.700 [0.691, 0.709] & 0.194 [0.178, 0.213] & 0.402 [0.389, 0.417] \\
C-PR ($\lambda=0.25$) & 0.422 [0.407, 0.438] & 0.430 [0.418, 0.444] & 0.157 [0.143, 0.170] \\
C-PR ($\lambda=0.5$) & 0.400 [0.385, 0.416] & 0.449 [0.437, 0.463] & 0.142 [0.128, 0.156] \\
C-PR ($\lambda=0.75$) & 0.377 [0.361, 0.393] & 0.464 [0.452, 0.478] & 0.129 [0.115, 0.143] \\
S-PR & 0.671 [0.661, 0.681] & 0.289 [0.272, 0.308] & 0.398 [0.386, 0.414] \\
\bottomrule
\end{tabular}
}
\end{table}


점수들은 자기가 걷는 화살표에 따라 갈려요. 송금만 걷는 두 점수인 AWP와 C-PR ($\lambda=0$)은 송금 가족에서 $0.736$($[0.728,0.744]$)과 $0.700$($[0.691,0.709]$)에 이르고, 허락 가족에서 $0.226$($[0.210,0.245]$)과 $0.194$($[0.178,0.213]$)에 이르러요. 다시 출발할 곳을 허락 걷기에서 얻는 S-PR은 $0.671$($[0.661,0.681]$)과 $0.289$($[0.272,0.308]$)에 이르러요. 허락만 걷는 두 점수인 엔도스랭크와 C-PR ($\lambda=1$)은 허락 가족에서 $0.560$($[0.548,0.571]$)과 $0.561$($[0.550,0.572]$)에 이르고, 송금 가족에서는 둘 다 $0.196$($[0.180,0.216]$)이에요. 그리고 C-PR의 중간 섞은 점수들은 두 층 사이에 있어요. 시빌 안정성 가족에서는 송금 위를 걷는 점수들이 앞서요. AWP는 $0.419$($[0.408,0.434]$)이고, 엔도스랭크는 $0.026$($[0.012,0.041]$)에 머물러요. 이 일치의 일부는 만든 방식 때문에 저절로 생겨요. 각 가족의 주된 대리 지표는 그래프 하나와 같은 이벤트에서 나와요. 입차수와 받은 양의 합은 송금에서 나오고, 받은 허락 수와 받은 허락 금액은 마지막 허락에서 나와요. 그리고 페이지랭크는 자기 그래프의 입차수를 매끄럽게 다듬어요. 그래서 점수가 자기 화살표 종류의 대리 지표와 맞는 것은 의도한 구성개념(점수가 재려고 하는 생각) 점검이에요. 이것을 외부 시험(바깥 자료로 하는 시험)으로 칠 수는 없어요. \ref{sec:holdout-results}절의 홀드아웃은 같은 질문을 기간 밖에서 물어요.

부록~\ref{sec:proxy-family-detail}는 엔도스랭크와 AWP의 지표 가족 평균을 대리 지표 하나하나로 나누어 보여 줘요. AWP의 계수 가운데 가장 큰 것은 송금 입차수와의 계수예요($\tau=0.742$, $[0.732,0.752]$). 받은 양의 합과의 계수도 거의 그만큼 높아요($0.730$, $[0.720,0.740]$). AWP의 무게를 보면 예상할 수 있는 결과예요. 기울기를 송금의 중앙값에 맞추었기 때문에(\ref{sub:awp-graph}절), 화살표 무게는 송금 횟수뿐 아니라 보낸 금액에 따라서도 커져요. 등록 분석에서는 원래 단위(소수점 조정을 하지 않은 가장 작은 단위)에서 $b=1$이었어요. 그래서 $V$는 거의 모든 송금을 대략 한 번으로 셌어요. 그리고 AWP는 받은 양의 합보다 입차수를 훨씬 가깝게 따라갔어요(부록~\ref{sec:map-registered}). 그렇다고 AWP가 입차수와 똑같아지지도 않아요. 퍼뜨림(전파), 시간 무게, 활발한 보낸 쪽에서 다시 출발하기가 두 계약 사이의 순서 가운데 눈에 띄는 몫을 바꿔 놓아요. 다만 기간 밖에서는 이렇게 바뀐 순서가 새 송금자를 입차수보다 못하게 줄 세워요(\ref{sec:holdout-results}절). 시빌 안정성 가족에서 AWP는 활동 기간($0.588$, $[0.575,0.600]$)과 활동한 달 수($0.616$, $[0.603,0.629]$)를 따라가요. 그리고 들어오는 상대 비율에서는 양수예요($0.055$, $[0.033,0.079]$).

엔도스랭크는 받은 허락 금액($0.469$, $[0.454,0.484]$)보다 받은 허락 수($0.650$, $[0.639,0.661]$)를 더 가깝게 따라가요. 엔도스랭크의 금액 무게는 허락 금액의 중앙값인 100달러의 두 배에서 $0.8$이에요. 그래서 몇백 달러의 허락과 무제한 허락은 무게가 거의 같아요. 이와 달리 받은 허락 금액은 허락 하나하나를 상한인 19,000달러까지 세요. 코호트 계약은 모두 주인 세 명 이상에게서 허락을 받았어요. 그래서 엔도스랭크는 동점(값이 같은 것)을 거의 남기지 않고, 엔도스랭크의 계수는 거의 모든 계약의 순서에 바탕을 두어요(\ref{sec:rank-divergence}절). C-PR ($\lambda=1$)은 상한으로 자른 금액으로 같은 허락 그래프를 걷고, 마찬가지로 고르게 다시 출발해요(균등 재시작). 이 점수는 모든 가족에서 엔도스랭크와의 차이가 $0.002$ 안에 머물러요.

주장~C4의 앞쪽 절반은 좁은 뜻에서 뒷받침돼요. 한 종류의 화살표 위를 걷는 점수는 모두 그 화살표 종류의 대리 지표와 가장 잘 맞고, 그 구간은 0에서 넉넉히 떨어져 있어요. 그리고 C-PR의 섞은 점수들은 모든 가족에서 자기의 두 층 사이에 있어요. 이 일치의 일부는 만든 방식 때문에 저절로 생겨요. \ref{sec:coupled-results}절에서는 기간 안에서 C-PR이 어디에 있는지 살펴봐요.

\dissertationSubheading[sec:rank-divergence]{순위 갈라짐(H1 / RQ1)}

가설~H1과 연구 질문 RQ1은 엔도스랭크와 AWP가 같은 계약들을 서로 다른 순서로 세우는지를 물어요. 두 점수 벡터 사이의 방법 간 켄달의 $\tau_b$는 $0.199$이고, 구간은 $[0.181,0.221]$예요(표~\ref{tab:tie-sensitivity}의 마지막 행). 그래서 이 구간에는 0도 1도 들어 있지 않아요. 두 순위는 서로 관계가 있지만, 서로 바꿔 쓸 수 있는 것과는 거리가 멀어요. 계약 쌍 $21{,}011{,}403$개 가운데 $96.4\%$는 두 점수 모두에서 순서가 정해져 있어요. 그리고 그 쌍들 가운데 $39.8\%$는 순서가 어긋나는 쌍이에요.

각 점수를 하나씩 따로 본 분포(주변 분포)는 그림~\ref{fig:score-distributions}에 있어요. 그리고 이 분포에서 동점은 거의 아무 역할도 하지 않아요. 엔도스랭크는 서로 다른 점수 $5{,}641$개를 주고, 계약 쌍의 $0.20\%$를 동점으로 둬요. AWP는 서로 다른 점수 $5{,}292$개를 주고, $3.38\%$를 동점으로 둬요. 두 점수 모두에서 가장 큰 동점 덩어리는 0점을 받은 계약들이에요. 엔도스랭크에서는 $197$개, AWP에서는 $1{,}192$개예요. 곧 각 그래프 밖에 있는 $191$개와 $1{,}192$개, 그리고 엔도스랭크에서는 금액이 너무 작아 금액 무게가 0과 다르지 않은 쌍만 가진 $6$개예요. 엔도스랭크에서는 계약 $63$개가 더, 들어오는 화살표가 없는 점의 점수 $c=1.93\times10^{-7}$을 함께 가져요. 그 가운데 $62$개는 허락 그래프에 주인으로만 들어 있어요. 나머지 하나는 1센트보다 적은 허락만 받았어요. 그 허락으로는 점수가 유효 숫자 아홉 자리까지 $c$ 그대로예요. 나머지 코호트 계약은 모두 이보다 큰 허락을 받았고, $c$보다 높은 점수를 받아요. 주변 분포로는 두 점수가 같은 계약들을 어떻게 줄 세우는지 보여 줄 수 없어요. 그래서 그림~\ref{fig:rank-scatter}은 무작위로 뽑은 계약 $3{,}000$개에서 두 순위를 서로 맞대어 그려요. 두 점수가 모두 위쪽에 두거나 모두 아래쪽에 두는 계약은 항등선(두 순위가 같은 곳을 이은 대각선) 가까이에 있어요. 그 사이에서는 점들이 사각형의 대부분에 흩어져 있어요.

0점 덩어리는 이 처리 과정의 규칙 하나에서 나와요. 어떤 그래프에 없는 계약은 0점을 받는다는 규칙이에요(\ref{sec:graph-construction}절). 분석 계획은 다른 방법, 곧 없는 계약을 외톨이 점(화살표가 하나도 없는 점)으로 남겨 두는 방법을 민감도 분석으로 정해 두었어요. 고르게 순간 이동(화살표를 무시하고 아무 점으로나 건너뛰기)을 하면, 외톨이 점은 들어오는 화살표가 없는 다른 점과 똑같은 점수를 받아요. 그래서 허락 그래프 밖에 있는 계약 $191$개는, 그 그래프의 점이 받을 수 있는 가장 낮은 점수인 $c$를 받아요. 그리고 순서의 맨 아래 가까이에서 $c$를 가진 계약 $63$개와 함께하게 돼요. 0 무게 쌍만 가진 계약 $6$개는 그래프의 점이지만, 풀이 프로그램이 빼는 점이에요. 그래서 이 계약들은 0점을 그대로 가져요. 방향이 바뀌는 쌍은 하나도 없어요. $191$개와 $63$개 사이의 쌍은 동점이 되고, $191$개와 $6$개 사이의 쌍은 순서가 생겨요. AWP에서 없는 계약 $1{,}192$개도 같은 방법으로 다루었어요. AWP의 다시 출발 규칙에서는, 아무것도 보내지 않은 외톨이 계약이 다시 출발 몫을 받지 못해요. 그래서 그 계약은 어느 쪽이든 0점이에요. 결과는 표~\ref{tab:tie-sensitivity}에 있어요. 어떤 계수도 소수 셋째 자리에서 바뀌지 않아요. 그리고 두 점수 사이의 일치는 $0.199$($[0.181,0.221]$) 그대로예요. 코호트는 받은 허락으로 정했어요. 그래서 거의 모든 코호트 계약은 허락 그래프의 스펜더이고, 없는 계약을 다루는 규칙은 거의 아무것도 바꾸지 않아요.

% Generated by experiments/scripts/export_latex.py -- do not edit by hand.
% Source: experiments/data/2-processed-tables-and-evaluations/revised-usd-12m/same-window/eval_summary.json
\begin{table}[htbp]
\centering
\caption[빠진 계약을 0점으로 매기거나 외톨이 점으로 둘 때의 같은 기간 지표 가족 평균.]{그래프에서 빠진 계약을 0점으로 매기거나 외톨이 점(화살표가 하나도 없는 점)으로 둘 때, 같은 기간의 지표 가족 평균 켄달의 $\tau$를 보여 줘요(매칭 코호트, 계약 $n=6{,}483$개). 주 분석에서는 어떤 방법의 그래프에 없는 계약을 0점으로 매겨요. 그 계약을 외톨이 점으로 두는 것은 분석 계획에 미리 정해 둔 민감도 분석이에요. 이 규칙의 영향을 받는 계약은 엔도스랭크에서 191개, AWP에서 1{,}192개예요(30개는 두 그래프 모두에 없어요). 엔도스랭크는 고르게 다시 출발하기 때문에, 외톨이 점은 들어오는 화살표가 없는 점과 같은 점수를 받아요. AWP는 송금을 보낸 주소에서만 다시 출발해요. 그래서 외톨이 계약은 점수 0을 그대로 갖고, AWP의 계수는 두 규칙에서 똑같아요. 마지막 행은 두 점수 사이의 일치예요. 신뢰구간(CI)은 계약을 다시 뽑아 구한 95\% 백분위 부트스트랩 구간이에요(짝지은 다시 뽑기 400번, 시드 42).}
\label{tab:tie-sensitivity}
\fitwidth{%
\begin{tabular}{lcccc}
\toprule
 & \multicolumn{2}{c}{빠진 계약을 0점으로} & \multicolumn{2}{c}{빠진 계약을 외톨이 점으로} \\
\cmidrule(lr){2-3} \cmidrule(lr){4-5}
지표 가족 & EndorseRank & AWP & EndorseRank & AWP \\
\midrule
송금 & 0.196 [0.180, 0.216] & 0.736 [0.728, 0.744] & 0.196 [0.180, 0.216] & 0.736 [0.728, 0.744] \\
허락 & 0.560 [0.548, 0.571] & 0.226 [0.210, 0.245] & 0.560 [0.548, 0.571] & 0.226 [0.210, 0.245] \\
시빌 안정성 & 0.026 [0.012, 0.041] & 0.419 [0.408, 0.434] & 0.026 [0.012, 0.041] & 0.419 [0.408, 0.434] \\
\midrule
EndorseRank 대 AWP & \multicolumn{2}{c}{0.199 [0.181, 0.221]} & \multicolumn{2}{c}{0.199 [0.181, 0.221]} \\
\bottomrule
\end{tabular}
}
\end{table}


\begin{figure}[htbp]
\centering
\includegraphics[width=0.9\textwidth]{\figpath{score-distributions.pdf}}
\caption[{매칭 코호트(계약 $n=6{,}483$개)에서 엔도스랭크와 AWP 점수의 주변 분포(점수가 양수인 계약, $\log_{10}$ 눈금).}]{매칭 코호트(계약 $n=6{,}483$개)에서 엔도스랭크와 AWP 점수의 주변 분포를 그렸어요. 점수가 양수인 계약을 $\log_{10}$ 눈금으로 그렸고, 개수도 로그 눈금이에요. 0점인 계약은 뺐어요. 이 계약들에는 그 그래프에서 무게가 양수인 화살표가 하나도 닿지 않아요. 엔도스랭크에서는 $197$개, AWP에서는 $1{,}192$개예요.}
\label{fig:score-distributions}
\end{figure}

\begin{figure}[htbp]
\centering
\includegraphics[width=0.8\textwidth]{\figpath{rank-scatter.pdf}}
\caption[{매칭 코호트에서 엔도스랭크와 AWP의 순위--순위 산점도(무작위로 뽑은 계약 $3{,}000$개).}]{매칭 코호트에서 엔도스랭크와 AWP의 순위를 서로 맞대어 찍은 산점도(점을 흩어 찍은 그림)예요(계약 $6{,}483$개 가운데 무작위로 뽑은 $3{,}000$개, 시드 42). 점수가 같으면 켄달의 $\tau_b$에서처럼 순위의 평균을 함께 가져요. 오른쪽 끝의 세로줄에는, 엔도스랭크에서 0점인 $197$개 가운데 뽑힌 계약들이 있어요. 맨 위의 가로줄에는, AWP에서 0점인 $1{,}192$개 가운데 뽑힌 계약들이 있어요. 점선으로 그린 항등선 위의 점은 두 순서에서 순위가 같아요. 범례(그림 속 설명 상자)에는 모든 계약에서 구한 $\tau_b$와 그 95\% 구간이 있어요.}
\label{fig:rank-scatter}
\end{figure}

H1은 뒷받침돼요. 두 점수는 코호트를 서로 다르게 줄 세워요. 그리고 그 차이는 동점에서도, 없는 계약을 다루는 규칙에서도 나오지 않아요. 송금에 바탕을 둔 점수에 기대는 서비스는, 허락에 바탕을 둔 점수로 바꾸거나 그것과 섞어도 계약의 순서가 그대로 남으리라고 가정할 수 없어요. 주장~C4의 뒤쪽 절반인 이 갈라짐은 결합 걷기의 전제(먼저 깔고 가는 생각)이기도 해요. 한 층 걷기는 저마다 한 가지 종류의 행동으로 계약을 줄 세워요. \ref{sec:cpr-results}절은 두 종류의 화살표를 함께 걷는 걷기 하나가 그 둘 가운데 어느 쪽보다도 나은지 물어요.

\dissertationSubheading[sec:robustness-results]{강건성 점검과 민감도 점검}

엔도스랭크와 AWP에 대해 세 가지 점검을 했어요. 세 점검은 결과가 감쇠 계수 값 하나, 금액 무게의 기울기 하나, 표본 정의 하나에 달려 있는지 물어요. 금액 기울기 점검은 계획에 있던 토큰 점검을 대신해요. 모든 그래프가 다섯 토큰을 남기게 되자, 토큰 점검은 할 까닭이 없어졌기 때문이에요(\ref{sub:top-tokens}절). 처음 두 점검은 표본을 그대로 둬요. 세 번째 점검은 표본 정의 자체를 바꿔요. 그래서 계수가 움직일 것으로 예상되고, 이 점검은 강건성의 증거가 아니에요. 네 번째 점검은 엔도스랭크가 다시 출발 규칙에 얼마나 기대는지 물어요.

\dissertationSubsubheading[sub:damping]{감쇠 계수}

표~\ref{tab:robustness-damping}에는 같은 그래프에서 $d\in\{0.75,0.85,0.95\}$일 때 지표 가족 평균을 다시 계산한 값이 있어요.

% Generated by experiments/scripts/export_latex.py -- do not edit by hand.
% Source: experiments/data/2-processed-tables-and-evaluations/revised-usd-12m/same-window/eval_summary.json
% Source: experiments/data/2-processed-tables-and-evaluations/revised-usd-12m/benchmark/benchmark.json
\begin{table}[htbp]
\centering
\caption[강건성 점검: 페이지랭크 감쇠 계수 $d$에 따른 지표 가족 평균 켄달의 $\tau$.]{강건성 점검(조건을 바꿔도 결과가 버티는지 보기)이에요. 매칭 코호트(계약 $n=6{,}483$개)에서 페이지랭크 감쇠 계수 $d \in \{0.75, 0.85, 0.95\}$를 쓸 때 엔도스랭크(ER)와 AWP의 지표 가족 평균 켄달의 $\tau$를 보여 줘요. 열 이름에서 허락은 허락 가족, 송금은 송금 가족, 시빌은 시빌 안정성 가족이에요. 실행 시간은 전체 그래프에서 시간을 잰 풀이 프로그램 호출 다섯 번의 평균이에요. $d=0.85$의 실행 시간은 표~\ref{tab:benchmark-runtime}에서 잰 값을 다시 쓴 것이고, 새로 재지 않았어요(모든 실행은 표~\ref{tab:benchmark-runs}에 있어요). 신뢰구간(CI)은 계약을 다시 뽑아 구한 95\% 백분위 부트스트랩 구간이에요(짝지은 다시 뽑기 400번, 시드 42).}
\label{tab:robustness-damping}
\fitwidth{%
\begin{tabular}{rccccccrr}
\toprule
$d$ & ER 허락 & AWP 허락 & ER 송금 & AWP 송금 & ER 시빌 & AWP 시빌 & ER (초) & AWP (초) \\
\midrule
0.75 & \begin{tabular}[c]{@{}c@{}}0.561\\[-1pt]{\footnotesize [0.549, 0.573]}\end{tabular} & \begin{tabular}[c]{@{}c@{}}0.229\\[-1pt]{\footnotesize [0.212, 0.248]}\end{tabular} & \begin{tabular}[c]{@{}c@{}}0.196\\[-1pt]{\footnotesize [0.180, 0.215]}\end{tabular} & \begin{tabular}[c]{@{}c@{}}0.740\\[-1pt]{\footnotesize [0.732, 0.748]}\end{tabular} & \begin{tabular}[c]{@{}c@{}}0.026\\[-1pt]{\footnotesize [0.012, 0.041]}\end{tabular} & \begin{tabular}[c]{@{}c@{}}0.427\\[-1pt]{\footnotesize [0.416, 0.441]}\end{tabular} & 1.986 & 5.956 \\
0.85 & \begin{tabular}[c]{@{}c@{}}0.560\\[-1pt]{\footnotesize [0.548, 0.571]}\end{tabular} & \begin{tabular}[c]{@{}c@{}}0.226\\[-1pt]{\footnotesize [0.210, 0.245]}\end{tabular} & \begin{tabular}[c]{@{}c@{}}0.196\\[-1pt]{\footnotesize [0.180, 0.216]}\end{tabular} & \begin{tabular}[c]{@{}c@{}}0.736\\[-1pt]{\footnotesize [0.728, 0.744]}\end{tabular} & \begin{tabular}[c]{@{}c@{}}0.026\\[-1pt]{\footnotesize [0.012, 0.041]}\end{tabular} & \begin{tabular}[c]{@{}c@{}}0.419\\[-1pt]{\footnotesize [0.408, 0.434]}\end{tabular} & 2.657 & 8.530 \\
0.95 & \begin{tabular}[c]{@{}c@{}}0.558\\[-1pt]{\footnotesize [0.547, 0.570]}\end{tabular} & \begin{tabular}[c]{@{}c@{}}0.217\\[-1pt]{\footnotesize [0.200, 0.235]}\end{tabular} & \begin{tabular}[c]{@{}c@{}}0.197\\[-1pt]{\footnotesize [0.181, 0.216]}\end{tabular} & \begin{tabular}[c]{@{}c@{}}0.723\\[-1pt]{\footnotesize [0.715, 0.731]}\end{tabular} & \begin{tabular}[c]{@{}c@{}}0.027\\[-1pt]{\footnotesize [0.013, 0.041]}\end{tabular} & \begin{tabular}[c]{@{}c@{}}0.406\\[-1pt]{\footnotesize [0.394, 0.420]}\end{tabular} & 8.113 & 20.805 \\
\bottomrule
\end{tabular}
}
\end{table}


엔도스랭크에서는 세 값 사이에서 어떤 지표 가족 평균도 $0.0031$보다 더 움직이지 않았어요. AWP에서는 $d$가 커질수록 송금 계수가 $0.740$($[0.732,0.748]$)에서 $0.723$($[0.715,0.731]$)으로 내려가요. 그리고 시빌 안정성 계수는 $0.427$($[0.416,0.441]$)에서 $0.406$($[0.394,0.420]$)으로 내려가요. 두 점수 어느 쪽의 모양도, 흔히 쓰는 $d=0.85$ 때문에 생긴 겉모습이 아니에요. 엔도스랭크가 잘 버티는 것은 대부분 그 구조 때문이에요. 주인들이 아무 허락도 받지 않은 스펜더에서는, 엔도스랭크가 $c\,(1+d\,\delta(v))$와 같아요. 여기서 $\delta(v)$는 주인 기준으로 나누어 맞춘(정규화한) 들어오는 무게의 합이에요(\ref{sub:theory-closed-form}절). 이런 스펜더들 사이에서 이 식이 주는 순서는 $d$에 달려 있지 않아요. $T_{\mathrm{obs}}$에서, 들어오는 화살표가 있는 코호트 계약 $6{,}224$개 가운데 이런 스펜더 $5{,}232$개 모두에서 이 식이 1\% 안으로 맞아요. 그래서 순위 상관은 거의 움직일 수 없어요. 실행 시간은 반복 횟수와 함께 늘고, 반복 횟수는 $d$와 함께 늘어요. 엔도스랭크는 $1.99$, $2.66$, $8.11$\,초가 걸려요(반복 $38$, $57$, $174$번). AWP는 $5.96$, $8.53$, $20.80$\,초가 걸려요($56$, $94$, $269$번). 시간 비율은 $3.00$, $3.21$, $2.56$이고, 이때 화살표 수 비율은 $2.21$이에요. 그림~\ref{fig:damping-sensitivity}에는 지표 가족 평균이 그려져 있어요.

\begin{figure}[htbp]
\centering
\includegraphics[width=0.85\textwidth]{\figpath{damping-sensitivity.pdf}}
\caption[{$d\in\{0.75,0.85,0.95\}$에서 지표 가족 평균 켄달의 $\tau$의 민감도(표~\ref{tab:robustness-damping}의 값).}]{엔도스랭크와 AWP의 지표 가족 평균 켄달의 $\tau$가 감쇠 계수에 따라 얼마나 바뀌는지 보여 줘요. 송금, 허락, 시빌 안정성 가족에서 $d\in\{0.75,0.85,0.95\}$에 걸쳐 그렸고, 95\% 부트스트랩 구간도 함께 그렸어요(값은 표~\ref{tab:robustness-damping}에서 가져왔어요. 매칭 코호트, 계약 $n=6{,}483$개).}
\label{fig:damping-sensitivity}
\end{figure}

\dissertationSubsubheading[sub:token-subgraph]{금액 기울기 민감도}

표~\ref{tab:robustness-slope}에는 금액 무게의 기울기 $b$를 10으로 나누거나 10을 곱해서(\ref{sub:top-tokens}절), $T_{\mathrm{obs}}$에서 엔도스랭크와 AWP를 다시 계산한 값이 있어요. 그래프와 대리 지표는 똑같이 두었어요.

% Generated by experiments/scripts/export_latex.py -- do not edit by hand.
% Source: experiments/data/2-processed-tables-and-evaluations/revised-usd-12m/robustness/robustness.json
% Source: experiments/data/0-usd-token-list-and-prices/usd-12m/constants.json
\begin{table}[htbp]
\centering
\caption[강건성 점검: 금액 기울기를 열 배 작게 하거나 크게 했을 때의 지표 가족 평균 켄달의 $\tau$.]{강건성 점검이에요. 금액 무게 $V$의 기울기 $b$를 열로 나누거나 열 배로 키웠을 때, 매칭 코호트(계약 $n=6{,}483$개)의 지표 가족 평균 켄달의 $\tau$를 주 분석 옆에 보여 줘요. 주 분석의 기울기는 금액이 중앙값(송금은 \$170, 허락은 \$100)일 때 $V=1/2$이 되게 정했어요. $b/10$은 그 점을 중앙값의 열 배로 옮기고, $10b$는 중앙값의 십분의 일로 옮겨요. 각 변형과, 같은 방법(엔도스랭크와 AWP)의 주 점수 사이의 순위 일치 $\tau$는 $b/10$에서 $0.979$와 $0.969$이고, $10b$에서 $0.971$과 $0.964$예요. 변형에 대해서는 구간을 계산하지 않았어요.}
\label{tab:robustness-slope}
\fitwidth{%
\begin{tabular}{lccccccccc}
\toprule
방법 & \multicolumn{3}{c}{송금} & \multicolumn{3}{c}{허락} & \multicolumn{3}{c}{시빌} \\
\cmidrule(lr){2-4} \cmidrule(lr){5-7} \cmidrule(lr){8-10}
 & $b/10$ & $b$ & $10b$ & $b/10$ & $b$ & $10b$ & $b/10$ & $b$ & $10b$ \\
\midrule
EndorseRank & 0.196 & 0.196 & 0.196 & 0.561 & 0.560 & 0.555 & 0.026 & 0.026 & 0.027 \\
AWP & 0.736 & 0.736 & 0.734 & 0.224 & 0.226 & 0.226 & 0.421 & 0.419 & 0.419 \\
\bottomrule
\end{tabular}
}
\end{table}


어느 기울기에서도 어떤 지표 가족 평균도 $0.0046$보다 더 움직이지 않았고, 모양도 그대로였어요. 허락 가족에서는 엔도스랭크가 앞서고, 송금 가족과 시빌 안정성 가족에서는 AWP가 앞서요. 더 가파른 기울기 $10b$에서는 무게가 2분의 1이 되는 금액이 중앙값 금액의 10분의 1이 되고, 화살표 무게는 이벤트를 센 수에 더 가까워져요. 이 기울기는 엔도스랭크의 허락 계수를 $0.560$에서 $0.555$로, AWP의 송금 계수를 $0.736$에서 $0.734$로 낮춰요. 더 완만한 기울기 $b/10$은 엔도스랭크의 계수를 $0.561$로 높이고, AWP의 계수는 $0.736$ 그대로 둬요. 각 변형은 자기 방법의 주 점수와 거의 같은 순서로 계약을 세워요($\tau$는 $0.964$에서 $0.979$까지). 변형들에 대해서는 구간을 계산하지 않았어요.

\dissertationSubsubheading[sub:sample-size]{표본 정의 민감도}

표~\ref{tab:robustness-sample-size}에는 크기에 따른 실행 시간을 볼 때 쓴 부분 표본에서 지표 가족 평균을 다시 계산한 값이 있어요.

% Generated by experiments/scripts/export_latex.py -- do not edit by hand.
% Source: experiments/data/2-processed-tables-and-evaluations/revised-usd-12m/robustness/robustness.json
% Source: experiments/data/2-processed-tables-and-evaluations/revised-usd-12m/benchmark/benchmark.json
\begin{table}[htbp]
\centering
\caption[평가 집합에 따른 지표 가족 평균의 민감도.]{평가 집합이 바뀌면 지표 가족 평균이 얼마나 움직이는지 보여 줘요. 매칭 코호트에서 SHA256으로 늘 같은 결과가 나오게 뽑은 부분 표본마다(시드 \texttt{contract-benchmark-v1}, 코호트 전체는 계약 $n=6{,}483$개) 엔도스랭크(ER)와 AWP의 지표 가족 평균 켄달의 $\tau$를 구했어요. 단계마다 그 단계의 계약에 닿는 화살표로 점수를 매겼어요. 열 이름에서 “허락”은 허락, “송금”은 송금, “시빌”은 시빌 안정성을 뜻해요. 행마다 평가한 계약 집합이 달라요. 그래서 이 $\tau$ 값들은 표~\ref{tab:alignment-hybrid}의 값과 비교할 수 없어요. 이 값들은 평가하는 집단에 따라 통계량이 얼마나 움직이는지 보여 줘요. 단계마다 구간은 계산하지 않았어요. 실행 시간은 표~\ref{tab:benchmark-scaling}의 측정값이에요.}
\label{tab:robustness-sample-size}
\fitwidth{%
\begin{tabular}{rccccccrrr}
\toprule
$n$ & ER 허락 & AWP 허락 & ER 송금 & AWP 송금 & ER 시빌 & AWP 시빌 & ER (초) & AWP (초) & AWP/ER 시간 \\
\midrule
1{,}000 & 0.665 & 0.228 & 0.225 & 0.752 & 0.049 & 0.441 & 0.498 & 1.307 & 2.6$\times$ \\
2{,}000 & 0.611 & 0.223 & 0.200 & 0.731 & 0.036 & 0.421 & 0.318 & 2.383 & 7.5$\times$ \\
5{,}000 & 0.577 & 0.229 & 0.201 & 0.735 & 0.033 & 0.420 & 1.672 & 6.620 & 4.0$\times$ \\
6{,}483 & 0.560 & 0.226 & 0.196 & 0.736 & 0.026 & 0.419 & 2.657 & 8.530 & 3.2$\times$ \\
\bottomrule
\end{tabular}
}
\end{table}


엔도스랭크가 가장 많이 움직여요. 엔도스랭크의 허락 계수는 $n=1{,}000$에서 $0.665$이고, 단계마다 내려가서 전체 코호트에서는 $0.560$이 돼요. 엔도스랭크의 송금 계수는 $0.225$에서 $0.196$으로 내려가요. AWP는 모든 가족에서 전체 코호트 값과의 차이가 $0.022$ 안에 머물러요. 단계들에 대해서는 구간을 계산하지 않았어요. 이런 변화 가운데 어느 것도 점수에 대한 증거가 아니에요. 한 단계는 뽑힌 계약에 닿는 화살표만 남겨요. 그래서 뽑히지 않은 코호트 계약을 지나는 길은 끊겨요. 그러면 각 점수는 뽑힌 계약 자신의 화살표를 센 수에 더 가까워지고, 그것이 바로 대리 지표가 세는 것이에요. 그래서 이 연구는 이 표를 표본 정의에 대한 민감도 분석으로 보고해요(\ref{sub:sample-size-protocol}절). 이 표의 값을 표~\ref{tab:alignment-hybrid}의 값과 견주어서는 안 돼요.

\dissertationSubsubheading[sub:endorserank-restarts]{엔도스랭크의 다시 출발}

엔도스랭크는 걷기를 고르게 다시 출발시켜요. 이와 달리 AWP는 각 주소가 최근에 보낸 활동에 비례해서 그 주소에서 걷기를 다시 출발시켜요(\ref{sub:endorserank-graph}절). 분석 계획은 엔도스랭크에 고르게 다시 출발하기를 정했어요. 그리고 그 변형을 그 옆에 함께 보고할 점수로 정했어요. 하지만 둘 사이의 비교는 하나도 정하지 않았어요. 표~\ref{tab:endorserank-restarts}에는 엔도스랭크의 두 버전의 점수가 있고, 그 옆에 AWP도 있어요.

% Generated by experiments/scripts/export_latex.py -- do not edit by hand.
% Source: experiments/data/2-processed-tables-and-evaluations/revised-usd-12m/same-window/eval_summary.json
% Source: experiments/data/2-processed-tables-and-evaluations/revised-usd-12m/holdout-w0/eval_summary.json
% Source: experiments/data/2-processed-tables-and-evaluations/revised-usd-12m/window-w1/eval_summary.json
\begin{table}[htbp]
\centering
\caption[AWP 옆에 놓은, 고르게 다시 출발하는 엔도스랭크와 활동 재시작을 쓰는 엔도스랭크.]{엔도스랭크를 고르게 다시 출발하게 한 것과, 허락하는 활동에 따라 무게를 단 활동 재시작을 쓰게 한 것을 AWP 옆에 놓았어요. 활동 재시작은 AWP의 규칙을 허락 그래프에 쓴 것이에요. 매칭 코호트(계약 $n=6{,}483$개)의 지표 가족 평균 켄달의 $\tau$를 보여 줘요. 또 두 스펜더 라벨과의 $\tau$를 W0에서(스펜더 계약 $n=7{,}914$개, 점수는 2026년 3월 31일에 동결), 그리고 W1에서($n=7{,}960$, 2026년 6월 30일에 동결) 보여 줘요. 두 엔도스랭크 행 모두 각 마지막 허락에 $\sigma(\Delta t)V(z)$만큼 무게를 줘요. 신뢰구간(CI)은 계약을 다시 뽑아 구한 95\% 백분위 부트스트랩 구간이에요(짝지은 다시 뽑기 400번, 시드 42).}
\label{tab:endorserank-restarts}
\fitwidth{%
\begin{tabular}{llccccccc}
\toprule
 & & \multicolumn{3}{c}{같은 기간, 매칭 코호트} & \multicolumn{2}{c}{W0, 스펜더} & \multicolumn{2}{c}{W1, 스펜더} \\
\cmidrule(lr){3-5} \cmidrule(lr){6-7} \cmidrule(lr){8-9}
방법 & 다시 출발 & 송금 & 허락 & 시빌 & 새 허락 쌍 & 새 송금자 & 새 허락 쌍 & 새 송금자 \\
\midrule
EndorseRank & 고르게 & \begin{tabular}[c]{@{}c@{}}0.196\\[-1pt]{\footnotesize [0.180, 0.216]}\end{tabular} & \begin{tabular}[c]{@{}c@{}}0.560\\[-1pt]{\footnotesize [0.548, 0.571]}\end{tabular} & \begin{tabular}[c]{@{}c@{}}0.026\\[-1pt]{\footnotesize [0.012, 0.041]}\end{tabular} & \begin{tabular}[c]{@{}c@{}}0.303\\[-1pt]{\footnotesize [0.287, 0.319]}\end{tabular} & \begin{tabular}[c]{@{}c@{}}0.272\\[-1pt]{\footnotesize [0.257, 0.286]}\end{tabular} & \begin{tabular}[c]{@{}c@{}}0.299\\[-1pt]{\footnotesize [0.284, 0.314]}\end{tabular} & \begin{tabular}[c]{@{}c@{}}0.259\\[-1pt]{\footnotesize [0.243, 0.275]}\end{tabular} \\
EndorseRank & 허락하는 활동에 따라 & \begin{tabular}[c]{@{}c@{}}0.251\\[-1pt]{\footnotesize [0.233, 0.271]}\end{tabular} & \begin{tabular}[c]{@{}c@{}}0.601\\[-1pt]{\footnotesize [0.590, 0.613]}\end{tabular} & \begin{tabular}[c]{@{}c@{}}0.056\\[-1pt]{\footnotesize [0.042, 0.072]}\end{tabular} & \begin{tabular}[c]{@{}c@{}}0.332\\[-1pt]{\footnotesize [0.315, 0.347]}\end{tabular} & \begin{tabular}[c]{@{}c@{}}0.270\\[-1pt]{\footnotesize [0.256, 0.285]}\end{tabular} & \begin{tabular}[c]{@{}c@{}}0.335\\[-1pt]{\footnotesize [0.319, 0.349]}\end{tabular} & \begin{tabular}[c]{@{}c@{}}0.236\\[-1pt]{\footnotesize [0.219, 0.251]}\end{tabular} \\
AWP & 송금하는 활동에 따라 & \begin{tabular}[c]{@{}c@{}}0.736\\[-1pt]{\footnotesize [0.728, 0.744]}\end{tabular} & \begin{tabular}[c]{@{}c@{}}0.226\\[-1pt]{\footnotesize [0.210, 0.245]}\end{tabular} & \begin{tabular}[c]{@{}c@{}}0.419\\[-1pt]{\footnotesize [0.408, 0.434]}\end{tabular} & \begin{tabular}[c]{@{}c@{}}0.196\\[-1pt]{\footnotesize [0.178, 0.213]}\end{tabular} & \begin{tabular}[c]{@{}c@{}}0.409\\[-1pt]{\footnotesize [0.396, 0.421]}\end{tabular} & \begin{tabular}[c]{@{}c@{}}0.157\\[-1pt]{\footnotesize [0.137, 0.178]}\end{tabular} & \begin{tabular}[c]{@{}c@{}}0.398\\[-1pt]{\footnotesize [0.386, 0.412]}\end{tabular} \\
\bottomrule
\end{tabular}
}
\end{table}


AWP의 다시 출발 규칙을 쓰면, 엔도스랭크의 같은 기간 허락 계수는 $0.560$($[0.548,0.571]$)에서 $0.601$($[0.590,0.613]$)로 올라가요. 송금 계수는 $0.196$($[0.180,0.216]$)에서 $0.251$($[0.233,0.271]$)로 올라가요. 이때 다시 출발 몫은 허락을 가장 많이 하는 주인들에게 가요. 그래서 이 점수는 스펜더의 주인이 몇 명인지뿐만 아니라, 그 주인들이 얼마나 활발한지로도 스펜더에 무게를 매겨요. 기간 밖에서 이 변형은 두 기간 모두에서 새 허락 쌍(새로 생긴 주인--스펜더 허락)을 더 잘 줄 세워요. 새 송금자(새로 송금을 보내온 주소)는 W0에서 더 잘 줄 세우지 못하고, W1에서는 더 못하게 줄 세워요. W0에서 이 변형은 새 허락 쌍에서 $0.332$($[0.315,0.347]$), 새 송금자에서 $0.270$($[0.256,0.285]$)에 이르러요. 고르게 다시 출발하는 버전과 견주면 $+0.029$($[0.023,0.037]$)와 $-0.002$($[-0.009,0.005]$)예요. W1에서는 $0.335$($[0.319,0.349]$)와 $0.236$($[0.219,0.251]$)에 이르고, 고르게 다시 출발하는 버전과 견주면 $+0.035$($[0.029,0.042]$)와 $-0.023$($[-0.032,-0.013]$)이에요. 표~\ref{tab:endorserank-restarts}에는 계수가 있고, 표~\ref{tab:fresh-posthoc}에는 W1의 대비(같은 다시 뽑은 표본에서 잰 두 값의 차이)가 있어요. 이 대비들은 라벨을 본 뒤에 계산했어요. 어느 버전도 새 허락 쌍을 받은 허락 수보다 더 잘 줄 세우지 못해요. W0에서 이 변형은 받은 허락 수보다 낮아요($-0.060$, $[-0.070,-0.052]$). 고르게 다시 출발하는 버전도 마찬가지예요($-0.090$, $[-0.100,-0.079]$; \ref{sec:holdout-results}절). \ref{sub:cpr-contrasts}절에서는 이 변형을 C-PR과 나란히 놓아 봐요.

\dissertationSubheading[sec:holdout-results]{한 층 걷기의 시간 홀드아웃(RQ3)}

같은 기간 점검으로는 점수와 그 점수 자신의 구성개념을 떼어 놓을 수 없어요. 그 대신 홀드아웃은, 어떤 기준 시점까지 잰 중심성(그래프에서 얼마나 중요한 자리에 있는지)이 그 뒤의 행동을 예측하는지(미리 맞히는지) 시험해요. \ref{sec:holdout-protocol}절에 따라, W0 스펜더 코호트의 점수를 $t_1$(2026년 3월 31일)에 동결해요. 그다음 2026년 4월 1일부터 6월 30일까지 계약마다 두 라벨을 세요. 하나는 보증(믿고 맡김) 구성개념인 새 주인--스펜더 허락 쌍이고, 다른 하나는 흐름 구성개념인 새 송금자예요. 표~\ref{tab:holdout-spenders}에서는 $t_1$에 계산한 모든 점수를 두 라벨과 맞대어 봐요. 두 원시 차수 기준선(비교 기준)도 함께 넣었어요. 표~\ref{tab:holdout-diff}에는 짝지은 대비가 있어요. 먼저 규칙~A의 두 대비가 있어요. 규칙~A는 이 연구의 로그를 하나도 뽑기 전에 분석 계획에서 정했어요. 그다음에는 라벨을 본 뒤에 계산한 대비들이 있어요. 이 절에서는 한 층 걷기들과 원시 차수들을 읽어요. 규칙~A와 C-PR, S-PR의 행은 \ref{sec:cpr-results}절에서 다뤄요.

% Generated by experiments/scripts/export_latex.py -- do not edit by hand.
% Source: experiments/data/2-processed-tables-and-evaluations/revised-usd-12m/holdout-w0/eval_summary.json
\begin{table}[htbp]
\centering
\caption[스펜더 코호트의 시간 홀드아웃 W0.]{스펜더 코호트(계약 $n=7{,}914$개)에서 해 본 시간 홀드아웃 W0이에요. 2026년 3월 31일에 동결한 점수와, 2026년 4월 1일부터 2026년 6월 30일까지 그 계약에서 센 두 라벨 사이의 켄달의 $\tau$를 보여 줘요. 두 라벨은 새로 생긴 주인--스펜더 허락 쌍의 수와, 그 계약에 처음으로 송금을 보낸 주소의 수예요. 코호트에는 동결일에 마지막 허락이 0보다 큰 스펜더 7{,}962개 가운데 계약 7{,}914개를 남겼어요. 코드가 없는 EOA 39개와 EIP-7702 위임이 있는 EOA 9개는 뺐어요. 계약 1{,}732개는 새 쌍을 하나 이상 받았고, 1{,}422개는 새 송금자를 하나 이상 받았어요. 기준선 행은 각 한 층 페이지랭크가 매끄럽게 다듬는 $t_1$의 원시 차수예요. C-PR ($\lambda=1$)과 C-PR ($\lambda=0$)은 결합 연산자의 허락 층과 송금 층을 각각 한 층만 걷는 것이에요. 신뢰구간(CI)은 계약을 다시 뽑아 구한 95\% 백분위 부트스트랩 구간이에요(짝지은 다시 뽑기 400번, 시드 42).}
\label{tab:holdout-spenders}
\fitwidth{%
\begin{tabular}{lcc}
\toprule
$t_1$의 점수 & 새 허락 쌍 ($\tau$ [95\% 신뢰구간]) & 새 송금자 ($\tau$ [95\% 신뢰구간]) \\
\midrule
$t_1$의 받은 허락 수 (기준선) & 0.392 [0.377, 0.406] & 0.284 [0.267, 0.299] \\
$t_1$의 송금 입차수 (기준선) & 0.203 [0.183, 0.224] & 0.447 [0.435, 0.459] \\
\midrule
EndorseRank & 0.303 [0.287, 0.319] & 0.272 [0.257, 0.286] \\
AWP & 0.196 [0.178, 0.213] & 0.409 [0.396, 0.421] \\
\midrule
C-PR ($\lambda=1$) & 0.298 [0.282, 0.313] & 0.267 [0.250, 0.282] \\
C-PR ($\lambda=0$) & 0.178 [0.160, 0.197] & 0.366 [0.351, 0.378] \\
C-PR ($\lambda=0.25$) & 0.329 [0.312, 0.343] & 0.345 [0.330, 0.357] \\
C-PR ($\lambda=0.5$) & 0.329 [0.312, 0.343] & 0.336 [0.322, 0.348] \\
C-PR ($\lambda=0.75$) & 0.327 [0.311, 0.342] & 0.327 [0.313, 0.339] \\
S-PR & 0.190 [0.171, 0.209] & 0.377 [0.363, 0.389] \\
\bottomrule
\end{tabular}
}
\end{table}


% Generated by experiments/scripts/export_latex.py -- do not edit by hand.
% Source: experiments/data/2-processed-tables-and-evaluations/revised-usd-12m/holdout-w0/eval_summary.json
\begin{table}[htbp]
\centering
\caption[스펜더 홀드아웃 W0에서 짝지은 부트스트랩으로 구한 차이.]{스펜더 홀드아웃 W0에서 짝지은 부트스트랩으로 구한 차이예요(계약 $n=7{,}914$개, $a$와 $b$가 함께 쓰는 다시 뽑은 표본 400개, 시드 42). 첫 번째 묶음에는 규칙 A의 주된 부분에 들어가는 두 대비가 있어요. 이 두 대비는 이 연구가 자료를 하나도 뽑기 전인 2026년 10월 9일에 커밋한 분석 계획에서 미리 정했고, 여기서는 고친 점수에 적용했어요. 하나는 새 허락 쌍에서 C-PR과 C-PR ($\lambda=1$)을 비교하고, 다른 하나는 새 송금자에서 C-PR과 C-PR ($\lambda=0$)을 비교해요. 두 번째 묶음은 라벨을 본 뒤에 계산했어요. 여기에는 다음이 들어 있어요. 허락 층만 걷는 C-PR ($\lambda=1$)이 $t_1$에서의 자기 원시 차수보다 더 얻은 값(증가분), S-PR에 대한 같은 비교들, 그리고 엔도스랭크와 AWP가 각자의 원시 차수보다 더 얻은 값이에요. C-PR의 다른 대비는 표~\ref{tab:cpr-contrasts}에 있어요. 규칙 A의 주된 부분은 C-PR이 새 허락 쌍에서 C-PR ($\lambda=1$)보다 더 나쁘지 않고, 새 송금자에서 C-PR ($\lambda=0$)보다 더 나아야 한다는 것이에요. “더 나쁘지 않음”은 성립하지만 “더 나음”은 성립하지 않아요. 그래서 이 규칙은 충족되지 않아요. 엔도스랭크와 AWP 사이의 차이는 표~\ref{tab:endorserank-awp-holdout}에 있어요.}
\label{tab:holdout-diff}
\fitwidth{%
\begin{tabular}{llccccc}
\toprule
대비 ($a$ 빼기 $b$) & 라벨 & $\tau(a)$ & $\tau(b)$ & $\Delta\tau$ & $\Delta\tau$의 95\% 신뢰구간 & $P(\Delta\tau>0)$ \\
\midrule
\multicolumn{7}{l}{\emph{규칙 A, 주된 부분, 계획에서 미리 정함}} \\
C-PR 빼기 C-PR ($\lambda=1$) & 새 허락 쌍 & 0.329 & 0.298 & +0.031 & [0.021, 0.042] & 1.000 \\
C-PR 빼기 C-PR ($\lambda=0$) & 새 송금자 & 0.336 & 0.366 & -0.030 & [-0.040, -0.020] & 0.000 \\
\midrule
\multicolumn{7}{l}{\emph{라벨을 본 뒤에 계산함}} \\
C-PR ($\lambda=1$) 빼기 $t_1$ 받은 허락 수 & 새 허락 쌍 & 0.298 & 0.392 & -0.094 & [-0.104, -0.083] & 0.000 \\
S-PR 빼기 C-PR ($\lambda=1$) & 새 허락 쌍 & 0.190 & 0.298 & -0.108 & [-0.128, -0.088] & 0.000 \\
S-PR 빼기 C-PR ($\lambda=0$) & 새 송금자 & 0.377 & 0.366 & +0.011 & [0.004, 0.018] & 1.000 \\
EndorseRank 빼기 $t_1$ 받은 허락 수 & 새 허락 쌍 & 0.303 & 0.392 & -0.090 & [-0.100, -0.079] & 0.000 \\
AWP 빼기 $t_1$ 입차수 & 새 송금자 & 0.409 & 0.447 & -0.038 & [-0.043, -0.032] & 0.000 \\
\bottomrule
\end{tabular}
}
\end{table}


첫 번째 발견은 한 층 걷기들에 대해 RQ3의 두 번째 부분에 답해요. 어느 걷기도, 다시 출발 규칙이 무엇이든, 자기 원시 차수보다 낫지 않아요. 활발한 보낸 쪽에서 다시 출발하는 AWP는 새 송금자를 $\tau=0.409$($[0.396,0.421]$)로 줄 세워요. 이것은 $t_1$의 송금 입차수($0.447$, $[0.435,0.459]$)보다 낮고, 짝지은 차이는 $-0.038$(구간 $[-0.043,-0.032]$)이에요. 고르게 다시 출발하며 걷는 송금 층은 더 낮아요. C-PR ($\lambda=0$)은 $0.366$($[0.351,0.378]$)에 이르러요(짝지은 대비는 계산하지 않았어요). 허락 쪽에서도 그냥 센 수가 이겨요. $t_1$의 받은 허락 수는 새 허락 쌍을 $\tau=0.392$($[0.377,0.406]$)로 예측해요. 엔도스랭크는 $0.303$($[0.287,0.319]$)에 이르고, 짝지은 차이는 $-0.090$($[-0.100,-0.079]$)이에요. C-PR ($\lambda=1$)은 $0.298$($[0.282,0.313]$)이고, 차이는 $-0.094$($[-0.104,-0.083]$)예요. 활동 재시작(활발한 주소에서 더 자주 다시 출발하기)을 써도 엔도스랭크는 그냥 센 수보다 낮아요(\ref{sub:endorserank-restarts}절). 석 달 기간 안에서 스펜더는 이미 가진 주인과 보낸 쪽이 많을수록 새 주인과 새 보낸 쪽을 더 많이 얻어요. 선호적 연결(이미 많이 가진 쪽이 더 많이 얻는 현상)이 예측하는 그대로예요\parencite{barabasi1999scaling}. 한 주소 주변만 보는 그냥 센 수는 이것을 바로 잡아내요. 엔도스랭크에서는 한 주인의 보증이, 그 주인이 허락한 모든 스펜더에게 나뉘어요(\ref{sub:theory-closed-form}절). 그래서 주인들이 다른 스펜더도 많이 허락한 스펜더는, 그냥 센 수만으로 놓일 자리보다 낮게 줄 세워져요. 이 라벨에서는 이렇게 나누는 것이 계수를 낮춰요. 이 차수 대비들은 라벨을 본 뒤에 계산했어요.

스펜더 코호트 전체를 놓고 읽으면, 각 라벨은 자기와 같은 종류의 화살표를 쓰는 점수가 가장 잘 줄 세워요. 새 허락 쌍에서는 받은 허락 수($0.392$, $[0.377,0.406]$)가 앞서요. 그다음은 C-PR 혼합 점수들($0.327$, $[0.311,0.342]$부터 $0.329$, $[0.312,0.343]$까지), 엔도스랭크, C-PR ($\lambda=1$)이에요. 반면에 송금 입차수($0.203$, $[0.183,0.224]$), AWP($0.196$, $[0.178,0.213]$), S-PR, C-PR ($\lambda=0$)은 더 낮아요. 새 송금자에서는 송금 입차수와 AWP가 앞서요. 허락 쪽 점수들은 더 낮고, 서로 가까워요. 엔도스랭크는 $0.272$($[0.257,0.286]$), C-PR ($\lambda=1$)은 $0.267$($[0.250,0.282]$), 받은 허락 수는 $0.284$($[0.267,0.299]$)에 이르고, 이들의 구간은 서로 겹쳐요. 엔도스랭크는 스펜더 $7{,}914$개에 서로 다른 점수 $7{,}368$개를 주고, 가장 큰 동점 덩어리에는 $116$개가 들어 있어요.

점수가 닿는 범위(커버리지)는 차수 비교를 설명하지 못해요. 하지만 새 허락 쌍에서 화살표 종류끼리 한 비교의 일부는 설명해요. 스펜더 7,914개 가운데 1,342개, 곧 여섯에 하나 정도는 동결 때까지 송금 화살표가 하나도 없어요. AWP, C-PR ($\lambda=0$), S-PR은 이 스펜더들에게 0점을 주고, 이들을 줄 세우는 것은 허락 화살표뿐이에요. 이 스펜더들은 나머지보다 새 허락 쌍을 더 자주 얻었고(26.7\% 대 20.9\%), 새 송금자는 거의 얻지 못했어요(1.4\% 대 21.3\%). 그래서 AWP의 0점들은 첫 번째 라벨에서는 AWP의 계수를 낮추고, 두 번째 라벨에서는 높여요. 표~\ref{tab:holdout-receiving}은 홀드아웃을, 동결 전에 다른 주소에게서 송금을 받은 적이 있는 스펜더 6,489개로 좁혀요. 거기서도 AWP는 새 송금자를 송금 입차수보다 못하게 줄 세워요($0.408$, $[0.393,0.422]$ 대 $0.459$, $[0.445,0.474]$; $-0.051$, $[-0.059,-0.043]$). 그리고 엔도스랭크는 새 허락 쌍에서 받은 허락 수보다 낮아요($0.299$, $[0.283,0.315]$ 대 $0.387$, $[0.370,0.402]$; $-0.087$, $[-0.096,-0.078]$). 새 허락 쌍에서 AWP는 이제 $0.328$($[0.312,0.343]$)에 이르러, 엔도스랭크보다 높고 받은 허락 수보다 낮아요(짝지은 대비는 계산하지 않았어요). 그래서 이 라벨에서 엔도스랭크가 AWP보다 앞선 것은 빠진 스펜더 1,425개에 기대요. 이 스펜더들에는 AWP가 0점을 주는 1,342개가 들어 있어요. 이렇게 좁히는 방법은 결과를 안 뒤에 골랐어요. 그 표의 C-PR 행은 \ref{sub:cpr-contrasts}절에서 읽어요.

% Generated by experiments/scripts/export_latex.py -- do not edit by hand.
% Source: experiments/data/2-processed-tables-and-evaluations/revised-usd-12m/holdout-w0/eval_summary.json
% Source: experiments/data/2-processed-tables-and-evaluations/revised-usd-12m/supplementary/supplementary.json
\begin{table}[htbp]
\centering
\caption[동결일 전에 송금을 받은 적이 있는 스펜더로 좁힌 스펜더 홀드아웃 W0.]{스펜더 홀드아웃 W0을 좁혀서 본 결과예요. 스펜더 계약 7{,}914개 가운데, 동결일 전에 다른 주소에서 송금을 받은 적이 있는 6{,}489개만 남겼어요($n=6{,}489$). 나머지 1{,}425개에는 다른 주소에서 들어오는 송금 화살표가 없어요. 그리고 스펜더 1{,}342개는 AWP 점수가 0이에요. 이렇게 좁히는 방법은 결과를 안 뒤에 골랐어요. 신뢰구간(CI)은 계약을 다시 뽑아 구한 95\% 백분위 부트스트랩 구간이에요(짝지은 다시 뽑기 400번, 시드 42). $^{a}$저장해 둔 점수로, 같은 다시 뽑은 표본을 써서 나중에 계산했어요. 이 부분 집합에서 C-PR을 AWP, 받은 허락 수와 견준 대비는 표~\ref{tab:cpr-contrasts}에 있어요.}
\label{tab:holdout-receiving}
\fitwidth{%
\begin{tabular}{lcc}
\toprule
$t_1$의 점수 & 새 허락 쌍 ($\tau$ [95\% 신뢰구간]) & 새 송금자 ($\tau$ [95\% 신뢰구간]) \\
\midrule
EndorseRank & 0.299 [0.283, 0.315] & 0.315 [0.300, 0.331] \\
AWP & 0.328 [0.312, 0.343] & 0.408 [0.393, 0.422] \\
C-PR ($\lambda=0.5$)$^{a}$ & 0.336 [0.321, 0.351] & 0.372 [0.359, 0.386] \\
\midrule
$t_1$의 받은 허락 수 & 0.387 [0.370, 0.402] & 0.335 [0.318, 0.350] \\
$t_1$의 송금 입차수 & 0.341 [0.325, 0.357] & 0.459 [0.445, 0.474] \\
\midrule
EndorseRank 빼기 $t_1$ 받은 허락 수 (새 허락 쌍) & \multicolumn{2}{c}{$\Delta\tau=-0.087$ [-0.096, -0.078]} \\
AWP 빼기 $t_1$ 입차수 (새 송금자) & \multicolumn{2}{c}{$\Delta\tau=-0.051$ [-0.059, -0.043]} \\
\bottomrule
\end{tabular}
}
\end{table}


표 전체에 두 가지 주의할 점이 있어요. 각 라벨은 같은 구성을 미래로 옮긴 것이에요(새 허락, 새 송금자). 그래서 이 스펜더 결과들 가운데 어느 것도 신용이나 손실과는 관계가 없어요. 아래의 트레이더 라벨은 두 종류의 화살표 어느 쪽으로도 만들지 않았지만, 신용 라벨도 아니에요. 또 두 라벨은 모두 듬성듬성해요. 스펜더 7,914개 가운데 1,732개가 새 허락 쌍을 얻었고, 1,422개가 새 송금자를 얻었어요. 그래서 모든 계수는 주로 이 작은 무리들이 어떻게 줄 세워지는지에 달려 있어요.

다른 후보 라벨들도 홀드아웃에서 계산했어요. 여기서 보고하지만, 뒷받침하는 증거로 내놓지는 않아요. 허락 취소 수에는 몇몇 토큰에서 끝까지 다 쓴 허락도 들어 있어요(\ref{sec:holdout-protocol}절). 허락 취소 수와 털어 가기 어림 규칙(drain heuristic, 토큰을 빼 가는 일을 대강 찾아내는 규칙)의 값은 둘 다 엔도스랭크와 함께 올라가고($\tau=0.343$, 구간 $[0.330,0.359]$, 그리고 $0.295$, 구간 $[0.279,0.312]$), 다른 모든 점수와도 함께 올라가요. 그리고 허락 취소는 원시 받은 허락 수와 함께 가장 가파르게 올라가요($0.384$, $[0.366,0.399]$). 보증으로 읽으면 부호가 틀렸어요. 믿을 만한 스펜더라면 허락을 더 자주가 아니라 덜 자주 잃어야 하기 때문이에요. 주인이 많은 스펜더는 주인들이 하는 모든 일을 더 많이 겪어요. 허락 취소와 털어 가기도 마찬가지예요. 그래서 이 라벨들은 믿음과 쓰임을 갈라낼 수 없어요.

\dissertationSubsubheading[sub:trader-results]{트레이더 결과(탐색적)}

\ref{sub:trader-holdout}절의 트레이더 실행은 W0의 동결을 GMX~V2 계약 트레이더(GMX V2에서 거래하는 계약)에 적용해요. 이 실행은 탐색적(미리 정하지 않고 살펴본) 결과로 보고해요. 이 실행의 청산(손실이 너무 커져 거래가 강제로 끝나는 일) 라벨은 \ref{sub:neutral-label-results}절의 비교를 위해 남겨 두어요. 표~\ref{tab:holdout-traders}에는 4월부터 6월까지 적어도 세 번 포지션을 닫은 계약 트레이더 87개에 대해, 이익 라벨 네 가지에서의 모든 점수가 있어요.

% Generated by experiments/scripts/export_latex.py -- do not edit by hand.
% Source: experiments/data/2-processed-tables-and-evaluations/revised-usd-12m/holdout-w0/eval_summary.json
\begin{table}[htbp]
\centering
\caption[GMX~V2 이익 라벨에 대한 W0 트레이더 분석.]{W0 트레이더 분석이에요. 2026년 3월 31일에 동결한 점수와, 2026년 4월 1일부터 2026년 6월 30일까지 센 GMX~V2 라벨 사이의 켄달의 $\tau$를 보여 줘요. 대상은 동결일 그래프에 있고 그 기간에 세 번 이상 닫은 GMX~V2 계약 계정 $n=87$개예요. 그 가운데 39개는 동결일까지 허락을 받은 적이 있었어요. 이 분석은 탐색적(미리 정하지 않고 살펴본)이에요. W0의 청산 라벨은 비교를 위해 남겨 두었어요(표~\ref{tab:neutral-label-contrasts}). 이익을 낸 닫기 수와 실현 이익은 닫은 횟수가 많을수록 커져요. 두 비율은 그렇지 않아요. 신뢰구간(CI)은 계약을 다시 뽑아 구한 95\% 백분위 부트스트랩 구간이에요(짝지은 다시 뽑기 400번, 시드 42).}
\label{tab:holdout-traders}
\fitwidth{%
\begin{tabular}{lcccc}
\toprule
$t_1$의 점수 & 이익을 낸 닫기 수 & 실현 이익 & 이익을 낸 비율 & 손실이 아닌 비율 \\
\midrule
$t_1$의 받은 허락 수 (기준선) & \begin{tabular}[c]{@{}c@{}}0.067\\[-1pt]{\footnotesize [-0.117, 0.225]}\end{tabular} & \begin{tabular}[c]{@{}c@{}}0.048\\[-1pt]{\footnotesize [-0.131, 0.222]}\end{tabular} & \begin{tabular}[c]{@{}c@{}}-0.160\\[-1pt]{\footnotesize [-0.338, 0.020]}\end{tabular} & \begin{tabular}[c]{@{}c@{}}-0.288\\[-1pt]{\footnotesize [-0.450, -0.126]}\end{tabular} \\
$t_1$의 송금 입차수 (기준선) & \begin{tabular}[c]{@{}c@{}}0.069\\[-1pt]{\footnotesize [-0.076, 0.210]}\end{tabular} & \begin{tabular}[c]{@{}c@{}}0.074\\[-1pt]{\footnotesize [-0.090, 0.229]}\end{tabular} & \begin{tabular}[c]{@{}c@{}}-0.117\\[-1pt]{\footnotesize [-0.275, 0.036]}\end{tabular} & \begin{tabular}[c]{@{}c@{}}-0.214\\[-1pt]{\footnotesize [-0.367, -0.062]}\end{tabular} \\
\midrule
EndorseRank & \begin{tabular}[c]{@{}c@{}}0.120\\[-1pt]{\footnotesize [-0.032, 0.260]}\end{tabular} & \begin{tabular}[c]{@{}c@{}}0.058\\[-1pt]{\footnotesize [-0.092, 0.196]}\end{tabular} & \begin{tabular}[c]{@{}c@{}}-0.047\\[-1pt]{\footnotesize [-0.193, 0.090]}\end{tabular} & \begin{tabular}[c]{@{}c@{}}-0.167\\[-1pt]{\footnotesize [-0.297, -0.025]}\end{tabular} \\
AWP & \begin{tabular}[c]{@{}c@{}}0.166\\[-1pt]{\footnotesize [0.015, 0.308]}\end{tabular} & \begin{tabular}[c]{@{}c@{}}0.197\\[-1pt]{\footnotesize [0.044, 0.332]}\end{tabular} & \begin{tabular}[c]{@{}c@{}}-0.121\\[-1pt]{\footnotesize [-0.264, 0.032]}\end{tabular} & \begin{tabular}[c]{@{}c@{}}-0.275\\[-1pt]{\footnotesize [-0.413, -0.145]}\end{tabular} \\
\midrule
C-PR ($\lambda=1$) & \begin{tabular}[c]{@{}c@{}}0.077\\[-1pt]{\footnotesize [-0.075, 0.233]}\end{tabular} & \begin{tabular}[c]{@{}c@{}}0.039\\[-1pt]{\footnotesize [-0.128, 0.197]}\end{tabular} & \begin{tabular}[c]{@{}c@{}}-0.076\\[-1pt]{\footnotesize [-0.240, 0.072]}\end{tabular} & \begin{tabular}[c]{@{}c@{}}-0.186\\[-1pt]{\footnotesize [-0.319, -0.045]}\end{tabular} \\
C-PR ($\lambda=0$) & \begin{tabular}[c]{@{}c@{}}0.166\\[-1pt]{\footnotesize [0.020, 0.300]}\end{tabular} & \begin{tabular}[c]{@{}c@{}}0.297\\[-1pt]{\footnotesize [0.133, 0.427]}\end{tabular} & \begin{tabular}[c]{@{}c@{}}-0.088\\[-1pt]{\footnotesize [-0.247, 0.062]}\end{tabular} & \begin{tabular}[c]{@{}c@{}}-0.233\\[-1pt]{\footnotesize [-0.375, -0.102]}\end{tabular} \\
C-PR ($\lambda=0.25$) & \begin{tabular}[c]{@{}c@{}}0.163\\[-1pt]{\footnotesize [0.012, 0.289]}\end{tabular} & \begin{tabular}[c]{@{}c@{}}0.289\\[-1pt]{\footnotesize [0.134, 0.414]}\end{tabular} & \begin{tabular}[c]{@{}c@{}}-0.082\\[-1pt]{\footnotesize [-0.234, 0.064]}\end{tabular} & \begin{tabular}[c]{@{}c@{}}-0.225\\[-1pt]{\footnotesize [-0.373, -0.089]}\end{tabular} \\
C-PR ($\lambda=0.5$) & \begin{tabular}[c]{@{}c@{}}0.172\\[-1pt]{\footnotesize [0.024, 0.300]}\end{tabular} & \begin{tabular}[c]{@{}c@{}}0.287\\[-1pt]{\footnotesize [0.131, 0.413]}\end{tabular} & \begin{tabular}[c]{@{}c@{}}-0.080\\[-1pt]{\footnotesize [-0.227, 0.062]}\end{tabular} & \begin{tabular}[c]{@{}c@{}}-0.222\\[-1pt]{\footnotesize [-0.365, -0.084]}\end{tabular} \\
C-PR ($\lambda=0.75$) & \begin{tabular}[c]{@{}c@{}}0.166\\[-1pt]{\footnotesize [0.021, 0.290]}\end{tabular} & \begin{tabular}[c]{@{}c@{}}0.273\\[-1pt]{\footnotesize [0.121, 0.397]}\end{tabular} & \begin{tabular}[c]{@{}c@{}}-0.076\\[-1pt]{\footnotesize [-0.228, 0.066]}\end{tabular} & \begin{tabular}[c]{@{}c@{}}-0.218\\[-1pt]{\footnotesize [-0.359, -0.081]}\end{tabular} \\
S-PR & \begin{tabular}[c]{@{}c@{}}0.168\\[-1pt]{\footnotesize [0.018, 0.301]}\end{tabular} & \begin{tabular}[c]{@{}c@{}}0.282\\[-1pt]{\footnotesize [0.118, 0.407]}\end{tabular} & \begin{tabular}[c]{@{}c@{}}-0.060\\[-1pt]{\footnotesize [-0.208, 0.092]}\end{tabular} & \begin{tabular}[c]{@{}c@{}}-0.203\\[-1pt]{\footnotesize [-0.354, -0.057]}\end{tabular} \\
\bottomrule
\end{tabular}
}
\end{table}


송금 화살표를 쓰는 걷기들은, 혼합 점수들까지 포함해서, 이익을 낸 닫기의 수와 실현 이익(실제로 얻은 이익)을 0보다 위에 있는 구간으로 줄 세워요. AWP는 $0.166$($[0.015,0.308]$)과 $0.197$($[0.044,0.332]$)에 이르러요. C-PR ($\lambda=0$), 세 혼합 점수, S-PR은 $0.163$($[0.012,0.289]$)부터 $0.172$($[0.024,0.300]$)까지, 그리고 $0.273$($[0.121,0.397]$)부터 $0.297$($[0.133,0.427]$)까지예요. 엔도스랭크, 활동 재시작을 쓴 그 변형, C-PR ($\lambda=1$), 두 원시 차수는 구간 안에 0이 들어 있어요. 그리고 송금 입차수는 두 라벨 모두에서 낮아요($0.069$, $[-0.076,0.210]$와 $0.074$, $[-0.090,0.229]$). 짝지은 대비가 두 쪽을 갈라 놓는 것은 실현 이익에서 C-PR의 두 층 사이뿐이에요($-0.257$, $[-0.447,-0.071]$). 이익을 낸 닫기에서는 갈라 놓지 못해요(엔도스랭크에서 AWP를 뺀 값 $-0.046$, $[-0.205,0.113]$; 표~\ref{tab:neutral-label-grid}). 두 라벨은 모두 계약이 얼마나 많이 거래하는지에 따라 커져요. 두 비율 라벨은 그렇지 않은데, 여기서는 모든 계수가 음수예요. 손실이 아닌 비율에서는 모든 구간이 0보다 아래에 있어요. 엔도스랭크의 $-0.167$($[-0.297,-0.025]$)부터 받은 허락 수의 $-0.288$($[-0.450,-0.126]$)까지예요. 이익을 낸 비율에서는 0보다 아래에 있는 구간이 하나도 없어요. 이 가운데 일부는 손실이 아닌 비율을 만든 방식 때문에 저절로 생겨요(\ref{sub:neutral-label-results}절). 어떤 점수도 계약이 얼마나 거래를 잘하는지는 줄 세우지 못해요. 등록 기간(계획을 먼저 등록해 둔 기간)에서는 C-PR에게 자기만의 트레이더 라벨이 주어져요.

\dissertationSubsubheading[sub:fresh-single]{등록 기간 W1의 한 층 걷기}

등록 기간 W1(\ref{sub:fresh-holdout}절)은 점수를 2026년 6월 30일인 $T_{\mathrm{obs}}$에 동결하고, 2026년 7월부터 9월까지 같은 두 스펜더 라벨을 세요. W1의 등록 문서와 규칙~B는 C-PR에 관한 것이고, \ref{sub:fresh-results}절에서 보고해요. 동결일에 스펜더 계약 7,960개가 양수인 마지막 허락을 가지고 있었어요. 그 가운데 1,495개는 라벨 기간에 새 허락 쌍을 얻었고, 1,371개는 새 송금자를 얻었어요. 표~\ref{tab:fresh-posthoc}에는 엔도스랭크와, 활동 재시작을 쓴 그 변형이 더해져 있어요. 등록 문서는 이 둘을 정하지 않아요. 이들의 행은 어떤 규칙에도 들어가지 않아요. 그리고 이 표의 대비는 라벨을 본 뒤에 계산했어요. 변형의 대비는 저장해 둔 W1 점수로 계산했어요.

이 기간에서도 한 층만 걷는 두 걷기 모두에서 W0의 모양이 그대로 나타나요. 동결일의 받은 허락 수는 이번에도 새 허락 쌍을 엔도스랭크보다 더 잘 줄 세워요($0.381$, $[0.365,0.396]$ 대 $0.299$, $[0.284,0.314]$; $-0.082$, $[-0.090,-0.073]$). 그리고 송금 입차수는 이번에도 새 송금자를 AWP보다 더 잘 줄 세워요($0.439$, $[0.426,0.454]$ 대 $0.398$, $[0.386,0.412]$; $-0.041$, $[-0.047,-0.034]$). 활동 재시작을 쓴 엔도스랭크는 W0에서처럼 여전히 새 허락 쌍을 차수보다 못하게 줄 세워요($-0.046$, $[-0.055,-0.038]$). W1 트레이더의 청산 없이 닫은 비율에서 엔도스랭크는 $0.297$($[0.174,0.418]$)에 이르러요.

% Generated by experiments/scripts/export_latex.py -- do not edit by hand.
% Source: experiments/data/2-processed-tables-and-evaluations/revised-usd-12m/window-w1/eval_summary.json
% Source: experiments/data/2-processed-tables-and-evaluations/revised-usd-12m/supplementary/supplementary.json
\begin{table}[htbp]
\centering
\caption[기간 W1에서 엔도스랭크와 활동 재시작을 쓰는 그 변형.]{기간 W1의 코호트와 라벨에서 엔도스랭크와, 활동 재시작을 쓰는 엔도스랭크 변형을 보여 줘요(스펜더 계약 $n=7{,}960$개, 트레이더 $n=70$). 이 행들도, 그 아래의 대비들도 규칙 B에 들지 않아요. 그리고 이 대비들은 등록 문서에 이름이 적혀 있지 않았어요. 활동 재시작을 쓰는 엔도스랭크의 대비는 저장해 둔 W1 점수로 나중에 계산했어요. AWP, C-PR, 차수 행은 표~\ref{tab:fresh-holdout}의 값을 다시 적은 것이고, 모든 대비는 같은 계약 다시 뽑은 표본을 써요. 엔도스랭크와 AWP 사이의 차이는 표~\ref{tab:endorserank-awp-holdout}에, C-PR의 차이는 표~\ref{tab:cpr-contrasts}에 있어요. 신뢰구간(CI)은 계약을 다시 뽑아 구한 95\% 백분위 부트스트랩 구간이에요(짝지은 다시 뽑기 400번, 시드 42).}
\label{tab:fresh-posthoc}
\fitwidth{%
\begin{tabular}{lccc}
\toprule
 & \multicolumn{2}{c}{스펜더 ($n=7{,}960$)} & 트레이더 ($n=70$) \\
\cmidrule(lr){2-3} \cmidrule(lr){4-4}
2026년 6월 30일의 점수 & 새 허락 쌍 & 새 송금자 & 청산 없이 닫은 비율 \\
\midrule
동결일의 받은 허락 수 (기준선) & 0.381 [0.365, 0.396] & 0.269 [0.250, 0.285] & 0.394 [0.261, 0.524] \\
동결일의 송금 입차수 (기준선) & 0.165 [0.143, 0.186] & 0.439 [0.426, 0.454] & 0.253 [0.116, 0.363] \\
\midrule
EndorseRank & 0.299 [0.284, 0.314] & 0.259 [0.243, 0.275] & 0.297 [0.174, 0.418] \\
EndorseRank, 활동 재시작(허락하는 활동) & 0.335 [0.319, 0.349] & 0.236 [0.219, 0.251] & 0.286 [0.160, 0.404] \\
AWP & 0.157 [0.137, 0.178] & 0.398 [0.386, 0.412] & 0.263 [0.140, 0.379] \\
C-PR ($\lambda=0.5$) & 0.316 [0.300, 0.330] & 0.324 [0.308, 0.337] & 0.282 [0.162, 0.397] \\
\midrule
EndorseRank 빼기 받은 허락 수, 새 허락 쌍 & \multicolumn{3}{l}{$\Delta\tau=-0.082$ [-0.090, -0.073]} \\
AWP 빼기 송금 입차수, 새 송금자 & \multicolumn{3}{l}{$\Delta\tau=-0.041$ [-0.047, -0.034]} \\
\midrule
EndorseRank (활동 재시작) 빼기 받은 허락 수, 새 허락 쌍 & \multicolumn{3}{l}{$\Delta\tau=-0.046$ [-0.055, -0.038]} \\
EndorseRank (활동 재시작) 빼기 EndorseRank, 새 허락 쌍 & \multicolumn{3}{l}{$\Delta\tau=+0.035$ [0.029, 0.042]} \\
EndorseRank (활동 재시작) 빼기 EndorseRank, 새 송금자 & \multicolumn{3}{l}{$\Delta\tau=-0.023$ [-0.032, -0.013]} \\
\bottomrule
\end{tabular}
}
\end{table}


\dissertationSubheading[sec:cpr-results]{결합 걷기 C-PR(RQ4, RQ5)}

C-PR은 이 논문의 방법이에요. C-PR은 달러로 잰 허락 층과 달러로 잰 송금 층 위를 걷는 무작위 걷기 하나예요. 이 걷기는 각 점에서 두 이동 행렬(전이 행렬) 줄을 무게 $\lambda$로 섞어요(\ref{sub:coupled-graph}절). 분석 계획은 판정 규칙 하나를 정했고, 그 대상은 C-PR이에요. 규칙~A는 결합 걷기가 한 층만 걷는 자기 층들을 하나하나 이기기를 요구해요(RQ4). W1의 등록 문서는 규칙~B를 더했어요. 규칙~B는 새 자료에서 C-PR이 자기 송금 층보다 나아지기를 요구해요(RQ5). 두 규칙은 모두 C-PR을 자기 층인 C-PR ($\lambda=1$), C-PR ($\lambda=0$)과 견줘요. 그 대비만이 결합만 바꾸고 다른 것은 하나도 바꾸지 않기 때문이에요. C-PR은 금액 무게와 고르게 다시 출발하기에서도 엔도스랭크, AWP와 달라요. 두 규칙은 모두 고친 점수에 재분석으로 적용해요. 이 절은 먼저 기간 안에서 C-PR이 어디에 있는지 살펴봐요. 그다음 W0에서 규칙~A를, W1에서 규칙~B를 보고해요. 마지막으로 C-PR을 이 연구의 다른 모든 점수와 견준 대비를 보고하는데, 이 대비들은 아무것도 판정하지 않아요.

\dissertationSubsubheading[sec:coupled-results]{기간 안의 결합 걷기}

\ref{sub:coupled-graph}절에서 혼합 점수 두 개를 정의했어요. C-PR은 각 점의 이동 행렬 줄에서, 달러로 잰 허락 층과 송금 층을 무게 $\lambda$로 섞어요. 그 빼 보기 실험인 S-PR은 송금 층을 그대로 쓰지만, 허락 걷기를 바탕으로 다시 출발해요. 표~\ref{tab:tau-diff-hybrid}는 이 기간에서 두 혼합 점수의 짝지은 대비를 보여 줘요. 먼저 \ref{sec:statistical-tests}절의 여덟 개는 한 층만 걷는 두 점수와 견줘요. 그다음 여섯 개는 엔도스랭크와 AWP와 견주고, 그 가운데 두 개는 나중에 저장해 둔 점수로 계산했어요. RQ4와 RQ5의 답은 기간 밖에서 정해요. 여기서는 한 층 점수들이 통과하도록 만들어진 점검 위에서, 혼합 점수들이 어디에 있는지만 살펴봐요.

% Generated by experiments/scripts/export_latex.py -- do not edit by hand.
% Source: experiments/data/2-processed-tables-and-evaluations/revised-usd-12m/same-window/eval_summary.json
% Source: experiments/data/2-processed-tables-and-evaluations/revised-usd-12m/supplementary/supplementary.json
\begin{table}[htbp]
\centering
\caption[C-PR과 S-PR을 두 층과, C-PR을 엔도스랭크 및 AWP와 견준 같은 기간 짝지은 대비.]{같은 기간에서 구한 짝지은 대비예요. 결합 연산자(C-PR, $\lambda=0.5$)와 보증 씨앗 걷기(S-PR)를, 이 둘을 만드는 두 층과 견주었어요. 두 층은 허락 층(C-PR ($\lambda=1$))과 송금 층(C-PR ($\lambda=0$))이에요. 또 C-PR을 엔도스랭크 및 AWP와 견주었어요(매칭 코호트, 계약 $n=6{,}483$개). 값은 지표 가족 평균 $\tau$이고, $a$와 $b$에는 같은 계약 다시 뽑은 표본을 써요. 층과의 대비는 규칙 A의 보조 부분 옆에 두는 더 엄격한 점검이에요. 규칙 A의 보조 부분은 C-PR의 점 추정값을 표~\ref{tab:alignment-hybrid}에서 더 나은 층의 구간과 비교해요. 규칙 A의 보조 부분은 분석 계획에 미리 정해 두었어요. 여기서는 그 규칙을 고친 점수에 썼어요. 그 결과 보조 부분의 조건은 채워지지 않아요. C-PR의 점 추정값은 송금 가족에서 C-PR ($\lambda=0$)의 구간 밖에 있고, 허락 가족에서 C-PR ($\lambda=1$)의 구간 밖에 있어요. 그리고 시빌 안정성 가족에서 C-PR은 두 층을 모두 넘지는 못해요. 구간이 0보다 아래이면, 표본을 뽑을 때 생기는 흔들림을 넘어서 $a$가 $b$보다 낮다는 뜻이에요. 구간이 0보다 위이면, 그 흔들림을 넘어서 더 높다는 뜻이에요. 신뢰구간(CI)은 계약을 다시 뽑아 구한 95\% 백분위 부트스트랩 구간이에요(짝지은 다시 뽑기 400번, 시드 42). $^{a}$저장해 둔 점수로, 같은 다시 뽑은 표본을 써서 나중에 계산했어요. 송금 가족에서는 C-PR을 엔도스랭크와, 허락 가족에서는 AWP와 견주었어요. 두 가족은 각각 그 한 층 걷기와 가장 덜 맞는 가족이에요.}
\label{tab:tau-diff-hybrid}
\fitwidth{%
\begin{tabular}{lccccc}
\toprule
대비 ($a$ 빼기 $b$) & $\tau(a)$ & $\tau(b)$ & $\Delta\tau$ & $\Delta\tau$의 95\% 신뢰구간 & $P(\Delta\tau>0)$ \\
\midrule
\multicolumn{6}{l}{\emph{한 층만 걷는 두 층과 비교}} \\
C-PR 송금 빼기 C-PR ($\lambda=0$) 송금 & 0.400 & 0.700 & -0.300 & [-0.318, -0.281] & 0.000 \\
C-PR 허락 빼기 C-PR ($\lambda=1$) 허락 & 0.449 & 0.561 & -0.112 & [-0.124, -0.100] & 0.000 \\
C-PR 시빌 안정성 빼기 C-PR ($\lambda=0$) 시빌 안정성 & 0.142 & 0.402 & -0.260 & [-0.278, -0.241] & 0.000 \\
C-PR 시빌 안정성 빼기 C-PR ($\lambda=1$) 시빌 안정성 & 0.142 & 0.028 & +0.115 & [0.104, 0.123] & 1.000 \\
S-PR 송금 빼기 C-PR ($\lambda=0$) 송금 & 0.671 & 0.700 & -0.029 & [-0.036, -0.023] & 0.000 \\
S-PR 허락 빼기 C-PR ($\lambda=1$) 허락 & 0.289 & 0.561 & -0.272 & [-0.290, -0.253] & 0.000 \\
S-PR 시빌 안정성 빼기 C-PR ($\lambda=0$) 시빌 안정성 & 0.398 & 0.402 & -0.004 & [-0.008, 0.000] & 0.037 \\
S-PR 시빌 안정성 빼기 C-PR ($\lambda=1$) 시빌 안정성 & 0.398 & 0.028 & +0.371 & [0.351, 0.390] & 1.000 \\
\midrule
\multicolumn{6}{l}{\emph{EndorseRank, AWP와 비교}} \\
C-PR 송금 빼기 AWP 송금 & 0.400 & 0.736 & -0.336 & [-0.354, -0.316] & 0.000 \\
C-PR 허락 빼기 EndorseRank 허락 & 0.449 & 0.560 & -0.111 & [-0.122, -0.099] & 0.000 \\
C-PR 시빌 안정성 빼기 AWP 시빌 안정성 & 0.142 & 0.419 & -0.277 & [-0.294, -0.260] & 0.000 \\
C-PR 시빌 안정성 빼기 EndorseRank 시빌 안정성 & 0.142 & 0.026 & +0.116 & [0.105, 0.125] & 1.000 \\
C-PR 송금 빼기 EndorseRank 송금$^{a}$ & 0.400 & 0.196 & +0.204 & [0.191, 0.214] & 1.000 \\
C-PR 허락 빼기 AWP 허락$^{a}$ & 0.449 & 0.226 & +0.223 & [0.206, 0.240] & 1.000 \\
\bottomrule
\end{tabular}
}
\end{table}


C-PR은 모든 가족에서 자기의 두 층 사이에 있고, 허락 층에 더 가까이 있어요. 주된 설정에서 C-PR은 허락 가족에서 $0.449$($[0.437,0.463]$)에 이르러, C-PR ($\lambda=1$)보다 $0.112$ 낮아요($[-0.124,-0.100]$). 송금 가족에서는 $0.400$($[0.385,0.416]$)에 이르러, C-PR ($\lambda=0$)보다 $0.300$ 낮아요($[-0.318,-0.281]$). $\lambda$가 $0.25$에서 $0.75$로 커지면, C-PR은 세 가족 모두에서 허락 층 쪽으로 움직여요. AWP는 송금 가족을 C-PR보다 더 잘 줄 세우고($-0.336$, $[-0.354,-0.316]$), 엔도스랭크는 허락 가족을 더 잘 줄 세워요($-0.111$, $[-0.122,-0.099]$). 각 한 층 걷기가 가장 덜 따라가는 가족에서는 C-PR이 그 걷기보다 앞서요. C-PR은 송금 가족을 엔도스랭크보다 $0.204$ 더 잘 줄 세우고($[0.191,0.214]$), 허락 가족을 AWP보다 $0.223$ 더 잘 줄 세워요($[0.206,0.240]$). 한 층에만 나가는 화살표가 있는 점은 $\lambda$가 얼마이든 그 층을 따라가요(식~\eqref{eq:coupled-transition}). $\lambda$는 두 층 모두에 나가는 화살표가 있는 점에서만 작동해요. 허락 층의 점은 거의 모두 나가는 화살표가 있는 주인이에요. 그래서 주인이 송금도 보냈고 추출에서 그 송금을 남겼다면, 그곳마다 허락 줄이 걷기에 들어가요. 시빌 안정성 가족에서 C-PR은 자기 허락 층보다 높고($+0.115$, $[0.104,0.123]$), 자기 송금 층보다 낮아요($-0.260$, $[-0.278,-0.241]$).

S-PR은 모든 가족에서, 그 가족에서 앞선 층보다 낮거나 그 층과 비슷해요. 송금 걷기를 허락 걷기를 바탕으로 다시 출발하게 하면, 송금 계수는 $0.671$로 떨어져요($[0.661,0.681]$; C-PR ($\lambda=0$)과 견주어 $\Delta\tau=-0.029$, $[-0.036,-0.023]$). 허락 계수는 $0.289$까지만 올라가요($[0.272,0.308]$; C-PR ($\lambda=1$)과 견주어 $\Delta\tau=-0.272$, $[-0.290,-0.253]$). 그리고 시빌 보정 안정성은 송금 층과 비슷한 채로 남아요($-0.004$, $[-0.008,0.000]$). 재시작 벡터는 순간 이동을 할 때마다 걷기가 어디에 내려앉는지를 바꿔요. 이 벡터는 허락 신호의 일부를 송금 걷기로 옮겨 오지만, 허락 화살표를 직접 걷는 것보다는 적게 옮겨 와요. 이것은 \ref{sub:cpr-rule-a}절이 기간 밖에서 얻는 발견을 같은 기간에서 보여 주는 짝이에요.

분석 계획에서 정한 대로, 규칙~A의 보조 부분은 다음을 요구해요. 송금 가족과 허락 가족에서는 C-PR의 점 추정값(구간이 아니라 값 하나로 낸 추정)이 더 나은 층의 95\% 구간 안에 있어야 해요. 그리고 시빌 보정 안정성에서는 C-PR이 두 층을 모두 넘어야 해요. $\lambda=0.5$에서 C-PR은 세 부분을 모두 통과하지 못해요. 등록 분석에서도 그랬어요(부록~\ref{sec:map-registered}). 송금 계수($0.400$)는 C-PR ($\lambda=0$)의 구간($[0.691,0.709]$)보다 아래에 있어요. 허락 계수($0.449$)는 C-PR ($\lambda=1$)의 구간($[0.550,0.572]$)보다 아래에 있어요. 그리고 시빌 보정 안정성에서 C-PR은 허락 층만 넘어요. 표~\ref{tab:tau-diff-hybrid}의 짝지은 대비도 같은 쪽을 가리켜요. 이 대비는 더 엄격한 점검으로서 규칙과 함께 보고해요. 같은 기간 표들에서 섞은 점수는 각 층이 자기 가족과 맞던 정도의 일부를 잃고, 어느 가족에서도 더 나은 층을 넘지 못해요. 남은 질문은 홀드아웃이 맡아요. 섞은 점수에 한 층 점수들에는 없는 예측 정보가 조금이라도 있는지 하는 질문이에요.

\dissertationSubsubheading[sub:cpr-rule-a]{기간 W0의 규칙~A}

결합 걷기는 기간 밖에서 규칙~A를 통과하지 못해요. 다만 탐색적 비교에서는 자기 송금 층만 걷는 것보다 나아요. 여기서 고친 점수에 재분석으로 적용한 그 규칙은, W0에서 C-PR에게 두 가지를 요구했어요. 새 허락 쌍에서는 C-PR ($\lambda=1$)보다 나쁘지 않아야 하고, 새 송금자에서는 C-PR ($\lambda=0$)을 이겨야 해요(표~\ref{tab:holdout-diff}의 첫 번째 묶음). $\lambda=0.5$에서 C-PR은 새 허락 쌍에서 허락 걷기보다 위에 있어요($\Delta\tau=+0.031$, 구간 $[0.021,0.042]$). 그래서 첫 번째 조건은 성립해요. 하지만 새 송금자에서는 송금 걷기보다 아래에 있어요($\Delta\tau=-0.030$, 구간 $[-0.040,-0.020]$). 그래서 두 번째 조건은 성립하지 않고, 그와 함께 주된 부분도 실패해요. 등록 분석에서도 그랬어요(표~\ref{tab:registered-vs-revised}). 규칙에 들어 있지 않은 비교로 송금 걷기 하나와 견주면, C-PR이 더 나아요. 새 허락 쌍에서 C-PR은 $0.329$($[0.312,0.343]$)에 이르고, 송금 걷기는 $0.178$($[0.160,0.197]$)이에요. 짝지은 차이는 $+0.151$이에요(같은 다시 뽑은 표본에서 구한 구간 $[0.130,0.170]$). 송금 걷기는 AWP처럼, 송금 화살표가 없는 스펜더에게 0점을 줘요. 그리고 나머지 스펜더들만 놓고 C-PR을 송금 걷기와 견주지는 않았어요. 그래서 이 이득의 일부는 그 스펜더들을 줄 세우는 데에서 나왔을 수 있어요. 이 대비는 라벨을 본 뒤에 계산했어요. 등록 분석에서는 같은 대비가 W1 등록 문서의 가설 F1이 되었어요(\ref{sub:fresh-results}절). 세 가지 $\lambda$ 값에 걸쳐 C-PR은 새 허락 쌍에서 거의 평평해요($\lambda=0.25$에서 $0.329$, $[0.312,0.343]$; $\lambda=0.75$에서 $0.327$, $[0.311,0.342]$). 그리고 새 송금자에서 C-PR의 점 추정값은 $\lambda$가 커질수록 $0.345$($[0.330,0.357]$)에서 $0.327$($[0.313,0.339]$)로 내려가요(짝지은 대비는 계산하지 않았고, $\lambda$를 맞추어 고르지도 않았어요).

씨앗을 쓰는 빼 보기 실험은 허락 신호가 어디에 있어야 하는지 보여 줘요. S-PR은 흐름 라벨을 송금 걷기보다 조금 더 잘 줄 세워요($\Delta\tau=+0.011$, 구간 $[0.004,0.018]$). 그리고 보증 라벨에서는 허락 걷기보다 낮아요($\Delta\tau=-0.108$, $[-0.128,-0.088]$). S-PR은 새 허락 쌍을 송금 걷기보다 아주 조금만 더 잘 줄 세워요($0.190$, $[0.171,0.209]$ 대 $0.178$). 걷기가 어디로 갈지를 송금 화살표가 정하면, 보증에서 끌어낸 재시작 분포(어디서 다시 출발하는지의 비율)는 별로 바꾸는 것이 없어요. 다만 S-PR도 송금 걷기처럼 송금 화살표가 없는 스펜더에게 0점을 줘요. 그래서 이 차이의 일부는 점수가 닿는 범위에서 나왔을 수 있어요. 스펜더 홀드아웃에서 허락 신호는, 허락 화살표가 걷기를 실어 나를 때 살아남아요. 허락 걷기에서는 혼자 다 실어 나르고, C-PR에서는 섞여서 실어 날라요. 하지만 허락 화살표가 걷기가 어디서 다시 출발할지만 정할 때는, 허락 신호의 대부분이 사라져요.

\dissertationSubsubheading[sub:fresh-results]{등록 기간 W1의 규칙~B}

\ref{sub:fresh-holdout}절의 등록 문서는 2026년 10월 9일 16:49:32~UTC에 공개 저장소에 커밋(저장소에 바뀐 내용을 기록함)했어요(커밋 \texttt{0d3a057}). 등록 분석의 W0 결과를 안 뒤였고, 이 연구가 7월--9월 로그를 뽑기 전이었어요. 등록 문서는 규칙~B와 허용 폭 $\delta=0.01$을 정했어요. 이 허용 폭은 등록 분석에서 F2 대비의 W0 구간의 반폭(구간 너비의 절반, $0.0040$)을 올림한 값이에요. 추출 매니페스트(실행 기록)에는 시험용 우회 설정을 쓴 기록이 없어요(부록~\ref{ch:appendix-eval}). 7월--9월 로그에는 계약 계정이 한 GMX~V2 닫기가 6,805개 있었고, 모두 해독되었어요. 고친 점수는 그 뒤에, 짝지은 다시 뽑기 400번과 시드 42로 평가했어요. 그리고 규칙~B와 R1--R3을 고친 점수에 재분석으로 적용해요. 등록 분석도 같은 판정에 이르렀지만, R1만은 달랐어요(부록~\ref{sec:map-registered}).

스펜더 코호트는 \ref{sub:fresh-single}절의 것이에요. 동결일 그래프에 있는 GMX~V2 계약 트레이더 70개가 라벨 기간에 적어도 세 번 포지션을 닫았어요. 이 가운데 9개는 청산이 적어도 한 번 있었고, 1개는 모든 닫기가 청산이었어요. 표~\ref{tab:fresh-holdout}에는 규칙~B의 모든 점수가 세 라벨에 대해 있고, 표~\ref{tab:fresh-contrasts}에는 대비와 그 판정이 있어요.

% Generated by experiments/scripts/export_latex.py -- do not edit by hand.
% Source: experiments/data/2-processed-tables-and-evaluations/revised-usd-12m/window-w1/eval_summary.json
\begin{table}[htbp]
\centering
\caption[기간 W1: 규칙 B의 세 라벨에 대한 켄달의 $\tau$.]{기간 W1의 결과예요. 2026년 6월 30일에 동결한 점수와 규칙 B의 세 라벨 사이의 켄달의 $\tau$를 보여 줘요. 라벨은 2026년 7월 1일부터 2026년 9월 30일까지 셌어요. 스펜더는 동결일에 마지막 허락이 0보다 큰 스펜더 계약 $n=7{,}960$개예요. 그 가운데 1{,}495개는 새 허락 쌍을 얻었고, 1{,}371개는 새 송금자를 얻었어요. 트레이더는 동결일 그래프에 있고 라벨 기간에 세 번 이상 닫은 GMX~V2 계약 계정 $n=70$개예요. 그 가운데 9개는 청산이 한 번 이상 있었고, 1개는 청산만 있었어요. C-PR ($\lambda=0$)은 규칙 B의 비교 대상이고, AWP는 규칙 B의 첫 번째 민감도 분석의 비교 대상이에요(표~\ref{tab:fresh-contrasts}). 엔도스랭크는 표~\ref{tab:fresh-posthoc}에 있어요. 신뢰구간(CI)은 계약을 다시 뽑아 구한 95\% 백분위 부트스트랩 구간이에요(짝지은 다시 뽑기 400번, 시드 42).}
\label{tab:fresh-holdout}
\fitwidth{%
\begin{tabular}{lccc}
\toprule
 & \multicolumn{2}{c}{스펜더 ($n=7{,}960$)} & 트레이더 ($n=70$) \\
\cmidrule(lr){2-3} \cmidrule(lr){4-4}
2026년 6월 30일의 점수 & 새 허락 쌍 & 새 송금자 & 청산 없이 닫은 비율 \\
\midrule
동결일의 받은 허락 수 (기준선) & 0.381 [0.365, 0.396] & 0.269 [0.250, 0.285] & 0.394 [0.261, 0.524] \\
동결일의 송금 입차수 (기준선) & 0.165 [0.143, 0.186] & 0.439 [0.426, 0.454] & 0.253 [0.116, 0.363] \\
\midrule
AWP & 0.157 [0.137, 0.178] & 0.398 [0.386, 0.412] & 0.263 [0.140, 0.379] \\
\midrule
C-PR ($\lambda=1$) & 0.295 [0.280, 0.310] & 0.257 [0.241, 0.274] & 0.302 [0.177, 0.425] \\
C-PR ($\lambda=0$) & 0.146 [0.126, 0.168] & 0.348 [0.335, 0.363] & 0.282 [0.162, 0.394] \\
C-PR ($\lambda=0.25$) & 0.315 [0.299, 0.328] & 0.332 [0.317, 0.346] & 0.282 [0.162, 0.397] \\
C-PR ($\lambda=0.5$) & 0.316 [0.300, 0.330] & 0.324 [0.308, 0.337] & 0.282 [0.162, 0.397] \\
C-PR ($\lambda=0.75$) & 0.316 [0.298, 0.330] & 0.315 [0.299, 0.329] & 0.284 [0.160, 0.398] \\
S-PR & 0.162 [0.140, 0.184] & 0.367 [0.354, 0.381] & 0.300 [0.175, 0.411] \\
\bottomrule
\end{tabular}
}
\end{table}


% Generated by experiments/scripts/export_latex.py -- do not edit by hand.
% Source: experiments/data/2-processed-tables-and-evaluations/revised-usd-12m/window-w1/eval_summary.json
\begin{table}[htbp]
\centering
\caption[기간 W1: 규칙 B의 대비와 판정, 그리고 민감도 분석.]{기간 W1에서 규칙 B의 짝지은 대비 F1--F6과 그 판정, 그리고 두 가지 민감도 분석을 보여 줘요. 규칙 B와 그 민감도 분석은 원래 단위로 한 분석을 위해 등록했어요. 여기서는 이것들을 고친 점수에 적용했어요. $\Delta\tau=\tau(a)-\tau(b)$이고, C-PR ($\lambda=0.5$)이 $a$예요. 우월성은 구간이 0보다 위에 있으면 성립해요. 비열등성은 구간의 아래 끝이 $-\delta=-0.01$보다 위에 있으면 성립해요. 규칙 A의 첫 부분인 F6a는 구간이 0을 품거나 0보다 위에 있으면 성립해요. 표~\ref{tab:fresh-holdout}에서처럼 F1, F2, F6a, F6b는 스펜더 계약($n=7{,}960$개)에서, F3--F5는 트레이더($n=70$)에서 구했어요. 신뢰구간(CI)은 계약을 다시 뽑아 구한 95\% 백분위 부트스트랩 구간이에요(짝지은 다시 뽑기 400번, 시드 42).}
\label{tab:fresh-contrasts}
\fitwidth{%
\begin{tabular}{lllcccll}
\toprule
 & 대비 & 라벨 & 검사 & $\tau(a)$ & $\tau(b)$ & $\Delta\tau$ [95\% 신뢰구간] & 판정 \\
\midrule
\multicolumn{8}{l}{\emph{규칙 B: C-PR ($\lambda=0.5$)과 그 송금 층만 걷는 C-PR ($\lambda=0$)의 비교}} \\
F1 & C-PR 대 C-PR ($\lambda=0$) & 새 허락 쌍 & 우월성 & 0.316 & 0.146 & +0.170 [0.149, 0.189] & 성립 \\
F2 & C-PR 대 C-PR ($\lambda=0$) & 새 송금자 & 비열등성, $-0.01$ & 0.324 & 0.348 & -0.025 [-0.038, -0.015] & 불성립 \\
F3 & C-PR 대 C-PR ($\lambda=0$) & 청산 없이 닫은 비율 & 비열등성, $-0.01$ & 0.282 & 0.282 & +0.000 [-0.010, 0.010] & 불성립 \\
\multicolumn{7}{l}{결정: F1, F2, F3 모두 성립해야 해요} & 충족 안 됨 \\
\midrule
\multicolumn{8}{l}{\emph{보고만 하고 결정에는 들어가지 않음}} \\
F4 & C-PR 대 C-PR ($\lambda=0$) & 청산 없이 닫은 비율 & 우월성 & 0.282 & 0.282 & +0.000 [-0.010, 0.010] & 불성립 \\
F5 & C-PR 대 C-PR ($\lambda=1$) & 청산 없이 닫은 비율 & 우월성 & 0.282 & 0.302 & -0.019 [-0.123, 0.089] & 불성립 \\
F6a & C-PR 대 C-PR ($\lambda=1$) & 새 허락 쌍 & 더 나쁘지 않음 & 0.316 & 0.295 & +0.021 [0.009, 0.032] & 성립 \\
F6b & C-PR 대 C-PR ($\lambda=0$) & 새 송금자 & 우월성 & 0.324 & 0.348 & -0.025 [-0.038, -0.015] & 불성립 \\
\midrule
\multicolumn{8}{l}{\emph{민감도 분석 1: AWP를 비교 대상으로 함}} \\
F1 & C-PR 대 AWP & 새 허락 쌍 & 우월성 & 0.316 & 0.157 & +0.159 [0.139, 0.178] & 성립 \\
F2 & C-PR 대 AWP & 새 송금자 & 비열등성, $-0.01$ & 0.324 & 0.398 & -0.075 [-0.085, -0.065] & 불성립 \\
F3 & C-PR 대 AWP & 청산 없이 닫은 비율 & 비열등성, $-0.01$ & 0.282 & 0.263 & +0.019 [-0.026, 0.064] & 불성립 \\
F4 & C-PR 대 AWP & 청산 없이 닫은 비율 & 우월성 & 0.282 & 0.263 & +0.019 [-0.026, 0.064] & 불성립 \\
\midrule
\multicolumn{8}{l}{\emph{민감도 분석 2: 그래프 밖의 계약을 외톨이 점으로 남김}} \\
F1 & C-PR 대 C-PR ($\lambda=0$) & 새 허락 쌍 & 우월성 & 0.316 & 0.148 & +0.168 [0.147, 0.187] & 성립 \\
F2 & C-PR 대 C-PR ($\lambda=0$) & 새 송금자 & 비열등성, $-0.01$ & 0.324 & 0.349 & -0.025 [-0.039, -0.016] & 불성립 \\
F3 & C-PR 대 C-PR ($\lambda=0$) & 청산 없이 닫은 비율 & 비열등성, $-0.01$ & 0.282 & 0.282 & +0.000 [-0.010, 0.010] & 불성립 \\
F4 & C-PR 대 C-PR ($\lambda=0$) & 청산 없이 닫은 비율 & 우월성 & 0.282 & 0.282 & +0.000 [-0.010, 0.010] & 불성립 \\
F5 & C-PR 대 C-PR ($\lambda=1$) & 청산 없이 닫은 비율 & 우월성 & 0.282 & 0.322 & -0.040 [-0.116, 0.035] & 불성립 \\
F6a & C-PR 대 C-PR ($\lambda=1$) & 새 허락 쌍 & 더 나쁘지 않음 & 0.316 & 0.295 & +0.021 [0.009, 0.032] & 성립 \\
F6b & C-PR 대 C-PR ($\lambda=0$) & 새 송금자 & 우월성 & 0.324 & 0.349 & -0.025 [-0.039, -0.016] & 불성립 \\
\bottomrule
\end{tabular}
}
\end{table}


규칙은 충족되지 않아요. F1은 성립해요. C-PR은 나중에 새 허락 쌍을 얻은 스펜더를 $\tau=0.316$($[0.300,0.330]$)으로 줄 세우고, 자기 송금 층만 걷는 걷기는 $0.146$($[0.126,0.168]$)이에요. 차이는 $+0.170$(구간 $[0.149,0.189]$)이에요. 그러니까 W0에서 본 허락 쪽의 이득이 다시 나타나요. 동결일의 받은 허락 수는 이 라벨을 $0.381$($[0.365,0.396]$)로 줄 세워, C-PR의 점 추정값보다 높아요. W0에서도 그랬어요($0.392$ 대 $0.329$). 그리고 두 기간 모두에서 C-PR은 받은 허락 수보다 낮아요(\ref{sub:cpr-contrasts}절). F1이 보여 주는 것은 송금 걷기보다 나은 이득이고, 그 일부는 송금 걷기가 0점을 주는 스펜더들에서 나왔을 수 있어요. F1은 허락한 주인 수(서로 다른 주인 수)보다 나은 이득을 보여 주지 않아요. F2는 성립하지 않아요. 새 송금자에서 C-PR은 $0.324$($[0.308,0.337]$)이고, 송금 걷기는 $0.348$($[0.335,0.363]$)이에요. 차이는 $-0.025$($[-0.038,-0.015]$)이고, 구간 전체가 허용 폭 $-0.01$보다 아래에 있어요. F3은 아슬아슬하게 성립하지 않아요. 트레이더의 청산 없이 닫은 비율에서 C-PR은 $0.282$($[0.162,0.397]$)에 이르고, 송금 걷기는 $0.282$($[0.162,0.394]$)에 이르러요. 차이는 $+0.000$($[-0.010,0.010]$)이고, 그 아래 끝이 $-0.01$보다 아주 조금 아래에 있어요. 규칙은 세 가설이 모두 성립해야 했어요. 그래서 등록 문서가 이런 경우에 대해 말한 대로, 이 논문은 C-PR이 송금만 걷는 걷기보다 나아진다고 주장하지 않아요.

등록 문서는 두 계수가 모두 0 가까이에 있어도 F3이 통과할 수 있다고 경고했어요. 여기서는 어느 쪽도 0 가까이에 있지 않아요. 표~\ref{tab:fresh-holdout}의 모든 점수는 이 라벨을 0보다 위에 있는 구간으로 줄 세워요. 동결일 송금 입차수의 $0.253$($[0.116,0.363]$)부터, 점 추정값이 가장 높은 받은 허락 수의 $0.394$($[0.261,0.524]$)까지예요. 걷기들 가운데 받은 허락 수와 견준 것은 엔도스랭크와 C-PR뿐이에요(\ref{sub:cpr-contrasts}절과 \ref{sub:neutral-label-results}절).

판정 옆에 함께 보고하는 결과 가운데 F4는 성립하지 않아요. 청산 없이 닫은 비율에서 C-PR은 송금 걷기보다 낫지 않아요. F5도 성립하지 않아요. 그 라벨에서 C-PR과 허락 걷기를 서로 가려낼 수 없기 때문이에요($-0.019$, $[-0.123,0.089]$). 규칙~A의 주된 부분인 F6도 이번에 다시 성립하지 않아요. C-PR은 새 허락 쌍에서 허락 걷기보다 위에 있어서($+0.021$, $[0.009,0.032]$) 나쁘지 않아야 한다는 조건은 만족해요. 하지만 새 송금자에서는 송금 걷기보다 아래에 있어요. W0에서처럼 C-PR의 계수는 새 허락 쌍에서는 $\lambda$에 걸쳐 거의 평평하고, 새 송금자에서는 $\lambda$가 커질수록 내려가요. S-PR은 새 송금자를 $0.367$($[0.354,0.381]$)로 줄 세워, 송금 걷기의 점 추정값보다 높아요. 그리고 새 허락 쌍을 $0.162$($[0.140,0.184]$)로 줄 세워, 허락 걷기($0.295$, $[0.280,0.310]$)보다 낮아요. W1에서 S-PR의 짝지은 대비는 하나도 등록하지 않았어요.

등록한 두 민감도 분석 어느 쪽도 판정을 바꾸지 않아요(표~\ref{tab:fresh-contrasts}). AWP를 비교 대상으로 삼아도 F1은 여전히 성립해요($+0.159$, $[0.139,0.178]$). 이 앞섬은 대부분 AWP가 0점을 주는 스펜더들에서 나오지만, 그것만은 아니에요(\ref{sub:cpr-contrasts}절). 그리고 F2는 뚜렷하게 성립하지 않아요. AWP는 새 송금자를 $0.398$($[0.386,0.412]$)로 줄 세우고, C-PR과 AWP의 차이는 $-0.075$($[-0.085,-0.065]$)예요. 이 구간은 모두 허용 폭보다 아래에 있어요. F3과 F4도 성립하지 않아요($+0.019$, $[-0.026,0.064]$). 트레이더 70개에서 이 구간의 너비는 $0.090$이에요. 그래서 F3은 약 $0.035$만큼 앞서야만 성립할 수 있었어요. F3이 성립하지 않았다는 것은 $0.026$만큼 뒤질 가능성을 남길 뿐이고, C-PR이 이 라벨을 더 못하게 줄 세운다는 것을 보여 주지는 않아요. 그래프 밖에 있는 계약을 외톨이 점으로 남겨 두면, 모든 판정이 그대로예요. 점 추정값은 많아야 $0.002$만큼 바뀌어요. F5만은 달라요. 이 규칙은 허락 그래프 밖에 있는 트레이더의 점수를 0에서, 들어오는 화살표가 없는 점의 점수로 올려 줘요. 그러면 청산 없이 닫은 비율에 대한 허락 걷기의 계수가 $0.302$에서 $0.322$로 올라가요. 그리고 C-PR과 허락 걷기의 차이는 $-0.040$($[-0.116,0.035]$)으로 벌어져요.

표~\ref{tab:fresh-traders}에는 판정에 들어가지 않는 W1 트레이더 라벨이 있어요. 구간은 \ref{sub:neutral-label-results}절의 다시 뽑은 표본 2,000개를 써요. 두 수 라벨(센 수로 된 라벨)에서는 이번에도 송금 화살표를 쓰는 걷기들이 앞서요. 반면에 송금 입차수는 두 라벨을 모두 0을 겨우 포함하는 구간으로 줄 세워요($0.156$, $[-0.012,0.306]$; $0.142$, $[-0.018,0.278]$). 그리고 이익을 낸 닫기에서 $\lambda=0.5$와 $0.75$인 C-PR의 구간, 실현 이익에서 세 혼합 점수와 S-PR의 구간에는 0이 들어 있어요. AWP는 이익을 낸 닫기($0.264$, $[0.103,0.408]$)와 실현 이익($0.280$, $[0.111,0.435]$)에서 점 추정값이 가장 높아요. 하지만 AWP가 C-PR ($\lambda=0$)보다 앞선다는 것은 어느 라벨에서도 가려지지 않아요(C-PR ($\lambda=0$)에서 AWP를 뺀 값 $-0.042$, $[-0.087,0.005]$와 $-0.052$, $[-0.108,0.002]$). 허락 쪽 점수들은 두 라벨 모두에서 0 가까이에 머물러요. 청산이 없었는지 나타내는 표시에서는 모든 점수가 양수이고, 받은 허락 수가 가장 높아요($0.407$, $[0.268,0.541]$). 두 비율에서는 받은 허락 수, C-PR ($\lambda=1$), 엔도스랭크, 활동 재시작을 쓴 그 변형의 구간이 0보다 아래에 있어요. 손실이 아닌 비율에서 받은 허락 수는 $-0.294$($[-0.468,-0.106]$)예요. 반면에 AWP와 C-PR ($\lambda=0$)은 양수이고, 그 구간에는 0이 들어 있어요.

% Generated by experiments/scripts/export_latex.py -- do not edit by hand.
% Source: experiments/data/2-processed-tables-and-evaluations/revised-usd-12m/window-w1/eval_summary.json
\begin{table}[htbp]
\centering
\caption[기간 W1: 다른 트레이더 라벨에 대한 켄달의 $\tau$.]{기간 W1의 트레이더예요. 2026년 6월 30일에 동결한 점수와, 2026년 7월 1일부터 2026년 9월 30일까지 센 다른 GMX~V2 라벨 사이의 켄달의 $\tau$를 보여 줘요. 대상은 그 기간에 포지션을 세 번 이상 닫은 계약 계정 $n=70$개예요. “청산 없음”은 그 계정이 닫은 포지션 가운데 청산이 하나도 없을 때 1이에요. 이익을 낸 닫기 수와 실현 이익은 닫은 횟수가 많을수록 커져요. 두 비율은 그렇지 않아요. 이 라벨들은 보고만 하고 규칙 B에는 들어가지 않아요. 구간은 표~\ref{tab:neutral-label-contrasts}에 쓴 다시 뽑은 표본을 써요. 신뢰구간(CI)은 계약을 다시 뽑아 구한 95\% 백분위 부트스트랩 구간이에요(짝지은 다시 뽑기 2{,}000번, 시드 42).}
\label{tab:fresh-traders}
\fitwidth{%
\begin{tabular}{lccccc}
\toprule
2026년 6월 30일의 점수 & 청산 없음 & 이익을 낸 닫기 수 & 실현 이익 & 이익을 낸 비율 & 손실이 아닌 비율 \\
\midrule
동결일의 받은 허락 수 (기준선) & \begin{tabular}[c]{@{}c@{}}0.407\\[-1pt]{\footnotesize [0.268, 0.541]}\end{tabular} & \begin{tabular}[c]{@{}c@{}}0.042\\[-1pt]{\footnotesize [-0.158, 0.243]}\end{tabular} & \begin{tabular}[c]{@{}c@{}}0.033\\[-1pt]{\footnotesize [-0.164, 0.230]}\end{tabular} & \begin{tabular}[c]{@{}c@{}}-0.221\\[-1pt]{\footnotesize [-0.406, -0.034]}\end{tabular} & \begin{tabular}[c]{@{}c@{}}-0.294\\[-1pt]{\footnotesize [-0.468, -0.106]}\end{tabular} \\
동결일의 송금 입차수 (기준선) & \begin{tabular}[c]{@{}c@{}}0.266\\[-1pt]{\footnotesize [0.125, 0.395]}\end{tabular} & \begin{tabular}[c]{@{}c@{}}0.156\\[-1pt]{\footnotesize [-0.012, 0.306]}\end{tabular} & \begin{tabular}[c]{@{}c@{}}0.142\\[-1pt]{\footnotesize [-0.018, 0.278]}\end{tabular} & \begin{tabular}[c]{@{}c@{}}-0.059\\[-1pt]{\footnotesize [-0.225, 0.099]}\end{tabular} & \begin{tabular}[c]{@{}c@{}}-0.071\\[-1pt]{\footnotesize [-0.225, 0.082]}\end{tabular} \\
\midrule
EndorseRank & \begin{tabular}[c]{@{}c@{}}0.303\\[-1pt]{\footnotesize [0.172, 0.422]}\end{tabular} & \begin{tabular}[c]{@{}c@{}}0.063\\[-1pt]{\footnotesize [-0.124, 0.253]}\end{tabular} & \begin{tabular}[c]{@{}c@{}}-0.011\\[-1pt]{\footnotesize [-0.211, 0.187]}\end{tabular} & \begin{tabular}[c]{@{}c@{}}-0.184\\[-1pt]{\footnotesize [-0.341, -0.014]}\end{tabular} & \begin{tabular}[c]{@{}c@{}}-0.265\\[-1pt]{\footnotesize [-0.411, -0.107]}\end{tabular} \\
EndorseRank, 활동 재시작(허락하는 활동) & \begin{tabular}[c]{@{}c@{}}0.292\\[-1pt]{\footnotesize [0.155, 0.415]}\end{tabular} & \begin{tabular}[c]{@{}c@{}}0.049\\[-1pt]{\footnotesize [-0.129, 0.237]}\end{tabular} & \begin{tabular}[c]{@{}c@{}}-0.019\\[-1pt]{\footnotesize [-0.195, 0.166]}\end{tabular} & \begin{tabular}[c]{@{}c@{}}-0.200\\[-1pt]{\footnotesize [-0.345, -0.037]}\end{tabular} & \begin{tabular}[c]{@{}c@{}}-0.275\\[-1pt]{\footnotesize [-0.418, -0.123]}\end{tabular} \\
AWP & \begin{tabular}[c]{@{}c@{}}0.265\\[-1pt]{\footnotesize [0.139, 0.383]}\end{tabular} & \begin{tabular}[c]{@{}c@{}}0.264\\[-1pt]{\footnotesize [0.103, 0.408]}\end{tabular} & \begin{tabular}[c]{@{}c@{}}0.280\\[-1pt]{\footnotesize [0.111, 0.435]}\end{tabular} & \begin{tabular}[c]{@{}c@{}}0.062\\[-1pt]{\footnotesize [-0.101, 0.213]}\end{tabular} & \begin{tabular}[c]{@{}c@{}}0.049\\[-1pt]{\footnotesize [-0.090, 0.185]}\end{tabular} \\
\midrule
C-PR ($\lambda=1$) & \begin{tabular}[c]{@{}c@{}}0.308\\[-1pt]{\footnotesize [0.175, 0.428]}\end{tabular} & \begin{tabular}[c]{@{}c@{}}0.009\\[-1pt]{\footnotesize [-0.167, 0.192]}\end{tabular} & \begin{tabular}[c]{@{}c@{}}-0.066\\[-1pt]{\footnotesize [-0.249, 0.120]}\end{tabular} & \begin{tabular}[c]{@{}c@{}}-0.208\\[-1pt]{\footnotesize [-0.370, -0.030]}\end{tabular} & \begin{tabular}[c]{@{}c@{}}-0.287\\[-1pt]{\footnotesize [-0.443, -0.121]}\end{tabular} \\
C-PR ($\lambda=0$) & \begin{tabular}[c]{@{}c@{}}0.286\\[-1pt]{\footnotesize [0.166, 0.400]}\end{tabular} & \begin{tabular}[c]{@{}c@{}}0.222\\[-1pt]{\footnotesize [0.072, 0.366]}\end{tabular} & \begin{tabular}[c]{@{}c@{}}0.228\\[-1pt]{\footnotesize [0.063, 0.388]}\end{tabular} & \begin{tabular}[c]{@{}c@{}}0.058\\[-1pt]{\footnotesize [-0.108, 0.210]}\end{tabular} & \begin{tabular}[c]{@{}c@{}}0.029\\[-1pt]{\footnotesize [-0.117, 0.170]}\end{tabular} \\
C-PR ($\lambda=0.25$) & \begin{tabular}[c]{@{}c@{}}0.284\\[-1pt]{\footnotesize [0.163, 0.399]}\end{tabular} & \begin{tabular}[c]{@{}c@{}}0.166\\[-1pt]{\footnotesize [0.011, 0.309]}\end{tabular} & \begin{tabular}[c]{@{}c@{}}0.168\\[-1pt]{\footnotesize [-0.002, 0.331]}\end{tabular} & \begin{tabular}[c]{@{}c@{}}0.008\\[-1pt]{\footnotesize [-0.157, 0.162]}\end{tabular} & \begin{tabular}[c]{@{}c@{}}-0.023\\[-1pt]{\footnotesize [-0.171, 0.117]}\end{tabular} \\
C-PR ($\lambda=0.5$) & \begin{tabular}[c]{@{}c@{}}0.283\\[-1pt]{\footnotesize [0.162, 0.398]}\end{tabular} & \begin{tabular}[c]{@{}c@{}}0.129\\[-1pt]{\footnotesize [-0.018, 0.271]}\end{tabular} & \begin{tabular}[c]{@{}c@{}}0.131\\[-1pt]{\footnotesize [-0.041, 0.297]}\end{tabular} & \begin{tabular}[c]{@{}c@{}}-0.025\\[-1pt]{\footnotesize [-0.185, 0.129]}\end{tabular} & \begin{tabular}[c]{@{}c@{}}-0.062\\[-1pt]{\footnotesize [-0.210, 0.080]}\end{tabular} \\
C-PR ($\lambda=0.75$) & \begin{tabular}[c]{@{}c@{}}0.284\\[-1pt]{\footnotesize [0.163, 0.398]}\end{tabular} & \begin{tabular}[c]{@{}c@{}}0.103\\[-1pt]{\footnotesize [-0.049, 0.254]}\end{tabular} & \begin{tabular}[c]{@{}c@{}}0.099\\[-1pt]{\footnotesize [-0.076, 0.267]}\end{tabular} & \begin{tabular}[c]{@{}c@{}}-0.051\\[-1pt]{\footnotesize [-0.212, 0.107]}\end{tabular} & \begin{tabular}[c]{@{}c@{}}-0.102\\[-1pt]{\footnotesize [-0.253, 0.047]}\end{tabular} \\
S-PR & \begin{tabular}[c]{@{}c@{}}0.304\\[-1pt]{\footnotesize [0.177, 0.418]}\end{tabular} & \begin{tabular}[c]{@{}c@{}}0.177\\[-1pt]{\footnotesize [0.033, 0.310]}\end{tabular} & \begin{tabular}[c]{@{}c@{}}0.146\\[-1pt]{\footnotesize [-0.031, 0.321]}\end{tabular} & \begin{tabular}[c]{@{}c@{}}0.009\\[-1pt]{\footnotesize [-0.143, 0.153]}\end{tabular} & \begin{tabular}[c]{@{}c@{}}-0.027\\[-1pt]{\footnotesize [-0.169, 0.111]}\end{tabular} \\
\bottomrule
\end{tabular}
}
\end{table}


\dissertationSubsubheading[sub:cpr-contrasts]{다른 모든 점수와 나란히 놓은 C-PR}

규칙~A와~B는 C-PR이 이 연구의 다른 점수들과 견주어 어디에 서는지에는 답하지 않아요. 표~\ref{tab:cpr-contrasts}은 C-PR을 그 점수들 하나하나와 두 원시 차수에 맞대어 봐요. 두 기간의 스펜더 라벨과 W1 트레이더의 청산 없이 닫은 비율에서 보고, 규칙에 쓴 다시 뽑은 표본을 그대로 써요. 표에 표시한 칸만 규칙에 들어요. 다른 대비는 모두 라벨을 본 뒤에 계산했고, 아무것도 판정하지 않아요.

% Generated by experiments/scripts/export_latex.py -- do not edit by hand.
% Source: experiments/data/2-processed-tables-and-evaluations/revised-usd-12m/holdout-w0/eval_summary.json
% Source: experiments/data/2-processed-tables-and-evaluations/revised-usd-12m/window-w1/eval_summary.json
% Source: experiments/data/2-processed-tables-and-evaluations/revised-usd-12m/supplementary/supplementary.json
\begin{table}[htbp]
\centering
\caption[기간 밖에서 C-PR을 다른 모든 점수, 그리고 원시 차수와 견준 대비.]{C-PR ($\lambda=0.5$)을 이 연구의 다른 모든 점수, 그리고 두 원시 차수와 짝지어 견준 대비예요. W0(계약 $n=7{,}914$개)과 W1($n=7{,}960$)의 스펜더 라벨, 그리고 W1 트레이더($n=70$)의 청산 없이 닫은 비율에서 구했어요. 칸마다 $\Delta\tau=\tau_{\text{C-PR}}-\tau_b$를 적고, 그 아래에 95\% 구간을 적었어요. 한 코호트 안에서는 모든 칸이 같은 다시 뽑은 표본을 써요. $^{\mathrm{A}}$를 붙인 칸은 분석 계획에 미리 정한 규칙 A의 대비예요. $^{\mathrm{B}}$는 W1의 등록 문서에 이름이 적힌 대비(F5와 F6a를 포함한 규칙 B)를, $^{\mathrm{S}}$는 규칙 B의 첫 번째 민감도 분석의 대비를 나타내요. 이 민감도 분석에서는 AWP가 비교 대상이에요. 이 규칙들은 원래 단위로 한 분석을 위해 정했고, 여기서는 고친 점수에 적용했어요. 나머지 칸은 모두 라벨을 본 뒤에 계산했고, 어떤 규칙에도 들지 않아요. 아래 묶음은 동결일 전에 다른 주소에서 송금을 받은 적이 있는 스펜더, 곧 AWP가 줄 세울 수 있는 스펜더에서 AWP, 받은 허락 수와의 대비를 다시 구한 것이에요. 이렇게 좁히는 방법도 결과를 안 뒤에 골랐어요. 신뢰구간(CI)은 계약을 다시 뽑아 구한 95\% 백분위 부트스트랩 구간이에요(짝지은 다시 뽑기 400번, 시드 42).}
\label{tab:cpr-contrasts}
\fitwidth{%
\begin{tabular}{lccccc}
\toprule
 & \multicolumn{2}{c}{W0 스펜더} & \multicolumn{2}{c}{W1 스펜더} & W1 트레이더 \\
\cmidrule(lr){2-3} \cmidrule(lr){4-5} \cmidrule(lr){6-6}
C-PR 빼기 & 새 허락 쌍 & 새 송금자 & 새 허락 쌍 & 새 송금자 & 청산 없이 닫은 비율 \\
\midrule
C-PR ($\lambda=1$) & \begin{tabular}[c]{@{}c@{}}+0.031$^{\mathrm{A}}$\\[-1pt]{\footnotesize [0.021, 0.042]}\end{tabular} & \begin{tabular}[c]{@{}c@{}}+0.069\\[-1pt]{\footnotesize [0.057, 0.080]}\end{tabular} & \begin{tabular}[c]{@{}c@{}}+0.021$^{\mathrm{B}}$\\[-1pt]{\footnotesize [0.009, 0.032]}\end{tabular} & \begin{tabular}[c]{@{}c@{}}+0.066\\[-1pt]{\footnotesize [0.054, 0.078]}\end{tabular} & \begin{tabular}[c]{@{}c@{}}-0.019$^{\mathrm{B}}$\\[-1pt]{\footnotesize [-0.123, 0.089]}\end{tabular} \\
C-PR ($\lambda=0$) & \begin{tabular}[c]{@{}c@{}}+0.151\\[-1pt]{\footnotesize [0.130, 0.170]}\end{tabular} & \begin{tabular}[c]{@{}c@{}}-0.030$^{\mathrm{A}}$\\[-1pt]{\footnotesize [-0.040, -0.020]}\end{tabular} & \begin{tabular}[c]{@{}c@{}}+0.170$^{\mathrm{B}}$\\[-1pt]{\footnotesize [0.149, 0.189]}\end{tabular} & \begin{tabular}[c]{@{}c@{}}-0.025$^{\mathrm{B}}$\\[-1pt]{\footnotesize [-0.038, -0.015]}\end{tabular} & \begin{tabular}[c]{@{}c@{}}+0.000$^{\mathrm{B}}$\\[-1pt]{\footnotesize [-0.010, 0.010]}\end{tabular} \\
EndorseRank & \begin{tabular}[c]{@{}c@{}}+0.026\\[-1pt]{\footnotesize [0.016, 0.037]}\end{tabular} & \begin{tabular}[c]{@{}c@{}}+0.064\\[-1pt]{\footnotesize [0.052, 0.075]}\end{tabular} & \begin{tabular}[c]{@{}c@{}}+0.017\\[-1pt]{\footnotesize [0.006, 0.029]}\end{tabular} & \begin{tabular}[c]{@{}c@{}}+0.064\\[-1pt]{\footnotesize [0.051, 0.076]}\end{tabular} & \begin{tabular}[c]{@{}c@{}}-0.015\\[-1pt]{\footnotesize [-0.116, 0.094]}\end{tabular} \\
EndorseRank (활동 재시작) & \begin{tabular}[c]{@{}c@{}}-0.003\\[-1pt]{\footnotesize [-0.014, 0.008]}\end{tabular} & \begin{tabular}[c]{@{}c@{}}+0.066\\[-1pt]{\footnotesize [0.055, 0.078]}\end{tabular} & \begin{tabular}[c]{@{}c@{}}-0.019\\[-1pt]{\footnotesize [-0.029, -0.008]}\end{tabular} & \begin{tabular}[c]{@{}c@{}}+0.087\\[-1pt]{\footnotesize [0.075, 0.102]}\end{tabular} & \begin{tabular}[c]{@{}c@{}}-0.004\\[-1pt]{\footnotesize [-0.102, 0.105]}\end{tabular} \\
AWP & \begin{tabular}[c]{@{}c@{}}+0.133\\[-1pt]{\footnotesize [0.114, 0.153]}\end{tabular} & \begin{tabular}[c]{@{}c@{}}-0.073\\[-1pt]{\footnotesize [-0.083, -0.063]}\end{tabular} & \begin{tabular}[c]{@{}c@{}}+0.159$^{\mathrm{S}}$\\[-1pt]{\footnotesize [0.139, 0.178]}\end{tabular} & \begin{tabular}[c]{@{}c@{}}-0.075$^{\mathrm{S}}$\\[-1pt]{\footnotesize [-0.085, -0.065]}\end{tabular} & \begin{tabular}[c]{@{}c@{}}+0.019$^{\mathrm{S}}$\\[-1pt]{\footnotesize [-0.026, 0.064]}\end{tabular} \\
S-PR & \begin{tabular}[c]{@{}c@{}}+0.139\\[-1pt]{\footnotesize [0.121, 0.156]}\end{tabular} & \begin{tabular}[c]{@{}c@{}}-0.041\\[-1pt]{\footnotesize [-0.049, -0.032]}\end{tabular} & \begin{tabular}[c]{@{}c@{}}+0.154\\[-1pt]{\footnotesize [0.134, 0.174]}\end{tabular} & \begin{tabular}[c]{@{}c@{}}-0.044\\[-1pt]{\footnotesize [-0.052, -0.037]}\end{tabular} & \begin{tabular}[c]{@{}c@{}}-0.017\\[-1pt]{\footnotesize [-0.094, 0.048]}\end{tabular} \\
동결일의 받은 허락 수 & \begin{tabular}[c]{@{}c@{}}-0.063\\[-1pt]{\footnotesize [-0.076, -0.052]}\end{tabular} & \begin{tabular}[c]{@{}c@{}}+0.052\\[-1pt]{\footnotesize [0.038, 0.065]}\end{tabular} & \begin{tabular}[c]{@{}c@{}}-0.065\\[-1pt]{\footnotesize [-0.077, -0.054]}\end{tabular} & \begin{tabular}[c]{@{}c@{}}+0.055\\[-1pt]{\footnotesize [0.040, 0.071]}\end{tabular} & \begin{tabular}[c]{@{}c@{}}-0.112\\[-1pt]{\footnotesize [-0.194, -0.026]}\end{tabular} \\
동결일의 송금 입차수 & \begin{tabular}[c]{@{}c@{}}+0.126\\[-1pt]{\footnotesize [0.104, 0.145]}\end{tabular} & \begin{tabular}[c]{@{}c@{}}-0.111\\[-1pt]{\footnotesize [-0.121, -0.100]}\end{tabular} & \begin{tabular}[c]{@{}c@{}}+0.151\\[-1pt]{\footnotesize [0.130, 0.173]}\end{tabular} & \begin{tabular}[c]{@{}c@{}}-0.115\\[-1pt]{\footnotesize [-0.128, -0.105]}\end{tabular} & \begin{tabular}[c]{@{}c@{}}+0.030\\[-1pt]{\footnotesize [-0.021, 0.099]}\end{tabular} \\
\midrule
\multicolumn{6}{l}{\emph{동결일 전에 송금을 받은 적이 있는 스펜더 (W0 $n=6{,}489$, W1 $n=6{,}567$)}} \\
AWP & \begin{tabular}[c]{@{}c@{}}+0.008\\[-1pt]{\footnotesize [-0.002, 0.020]}\end{tabular} & \begin{tabular}[c]{@{}c@{}}-0.036\\[-1pt]{\footnotesize [-0.044, -0.026]}\end{tabular} & \begin{tabular}[c]{@{}c@{}}+0.022\\[-1pt]{\footnotesize [0.011, 0.033]}\end{tabular} & \begin{tabular}[c]{@{}c@{}}-0.037\\[-1pt]{\footnotesize [-0.046, -0.028]}\end{tabular} & --- \\
동결일의 받은 허락 수 & \begin{tabular}[c]{@{}c@{}}-0.051\\[-1pt]{\footnotesize [-0.064, -0.037]}\end{tabular} & --- & \begin{tabular}[c]{@{}c@{}}-0.053\\[-1pt]{\footnotesize [-0.066, -0.041]}\end{tabular} & --- & --- \\
\bottomrule
\end{tabular}
}
\end{table}


허락만 걷는 엔도스랭크와 견주면, C-PR은 두 기간 모두에서 두 스펜더 라벨을 더 잘 줄 세워요. 새 허락 쌍에서는 W0에서 $0.026$($[0.016,0.037]$), W1에서 $0.017$($[0.006,0.029]$)만큼 앞서요. 새 송금자에서는 $0.064$($[0.052,0.075]$)와 $0.064$($[0.051,0.076]$)만큼 앞서요. W1 트레이더에서는 C-PR과 엔도스랭크를 서로 가려낼 수 없어요($-0.015$, $[-0.116,0.094]$). 활동 재시작을 쓴 엔도스랭크 변형은 W0의 새 허락 쌍에서 C-PR과 비슷하고($-0.003$, $[-0.014,0.008]$), W1에서는 새 허락 쌍을 C-PR보다 더 잘 줄 세워요($-0.019$, $[-0.029,-0.008]$). 새 송금자에서는 두 기간 모두 C-PR이 이 변형보다 앞서요($+0.066$, $[0.055,0.078]$와 $+0.087$, $[0.075,0.102]$). 기간 안에서 C-PR은 송금 가족과 시빌 안정성 가족에서 이 변형보다 앞서고($0.400$과 $0.142$ 대 $0.251$과 $0.056$), 허락 가족에서는 뒤져요($0.449$ 대 $0.601$; 표~\ref{tab:alignment-hybrid}과~\ref{tab:endorserank-restarts}).

송금만 걷는 AWP와 견주면, C-PR은 새 허락 쌍을 더 잘 줄 세우고(W0에서 $+0.133$, $[0.114,0.153]$, W1에서 $+0.159$, $[0.139,0.178]$), 새 송금자를 더 못하게 줄 세워요($-0.073$, $[-0.083,-0.063]$와 $-0.075$, $[-0.085,-0.065]$). 허락에서 앞서는 것은 대부분 AWP가 줄 세우지 못하는 스펜더들에서 나와요. AWP는 동결 때까지 송금 화살표가 없는 스펜더에게 모두 0점을 줘요. W0에서 이런 스펜더는 1,342개예요. 그리고 이 스펜더들은 나머지보다 새 허락 쌍을 더 자주 얻어요(\ref{sec:holdout-results}절). 동결 전에 다른 주소에게서 송금을 받은 적이 있는 스펜더들 사이에서는, W0의 새 허락 쌍에서 C-PR과 AWP가 비슷하고($+0.008$, $[-0.002,0.020]$), W1에서는 C-PR이 앞서요($+0.022$, $[0.011,0.033]$). 반면에 새 송금자는 두 기간 모두 AWP가 더 잘 줄 세워요($-0.036$, $[-0.044,-0.026]$와 $-0.037$, $[-0.046,-0.028]$)(표~\ref{tab:cpr-contrasts}의 아래 묶음. 이렇게 좁히는 방법도 결과를 안 뒤에 골랐어요). 기간 밖에서 C-PR이 송금 걷기에 보태는 것은, 주로 이 걷기가 동점으로 남기는 계약들 사이의 순서예요. 그 대가로, 이 걷기가 줄 세울 수 있는 계약들 사이에서는 나중에 들어오는 흐름의 순서가 더 나빠져요.

원시 차수와 견주면, C-PR은 어느 라벨도 같은 종류 화살표를 센 수보다 더 잘 줄 세우지 못해요. 새 허락 쌍에서는 두 기간 모두 받은 허락 수보다 낮아요($-0.063$, $[-0.076,-0.052]$와 $-0.065$, $[-0.077,-0.054]$). 새 송금자에서는 두 기간 모두 송금 입차수보다 낮아요($-0.111$, $[-0.121,-0.100]$와 $-0.115$, $[-0.128,-0.105]$). 자기의 빼 보기 실험인 S-PR과 견주면, C-PR은 새 허락 쌍을 더 잘 줄 세우고($+0.139$, $[0.121,0.156]$와 $+0.154$, $[0.134,0.174]$), 새 송금자를 더 못하게 줄 세워요($-0.041$, $[-0.049,-0.032]$와 $-0.044$, $[-0.052,-0.037]$). 트레이더 70개의 청산 없이 닫은 비율에서는 C-PR이 받은 허락 수보다 낮아요($-0.112$, $[-0.194,-0.026]$). C-PR의 다른 대비는 하나도 0을 벗어나지 않아요.

모두 합쳐 보면, C-PR은 두 기간 모두에서 모든 스펜더 라벨을 엔도스랭크보다 더 잘 줄 세워요. 그리고 나중의 허락을 AWP보다 더 잘 줄 세우는데, 대부분 AWP가 동점으로 남기는 계약들 덕분이에요. C-PR은 두 기간 모두에서, 각 화살표 종류의 라벨에서 그 종류를 그냥 센 수보다 낮아요. 나중의 송금자에서는 AWP보다 낮고, W1의 나중 허락에서는 활동 재시작을 쓴 엔도스랭크보다 낮아요. 그리고 트레이더의 청산 없이 닫은 비율에서는 받은 허락 수보다 낮아요. 이 대비들 가운데 어느 것도 미리 정해 두지 않았어요.

\dissertationSubheading[sub:neutral-label-results]{두 종류의 화살표 어느 쪽으로도 만들지 않은 결과 값}

GMX 트레이더 라벨은 이 자료에서 두 종류의 화살표 어느 쪽으로도 만들지 않은 유일한 결과 값이에요. \ref{sub:neutral-label-protocol}절의 비교는 규칙~B와 함께 등록했어요. 그때는 이 연구가 7월--9월 로그를 뽑기 전이었고, W0의 청산 라벨에서 점수를 평가하기 전이었어요. 여기서는 그 비교의 대비, 규칙, 문구를 고친 점수에 재분석으로 적용해요. W1이 주된 기간이고, W0이 보조 기간이에요. 구간은 짝지은 다시 뽑기 2,000번에서 나왔어요. 97.5\%라고 표시하지 않은 구간은 95\% 구간이에요. 97.5\% 구간은 판정 대비 P와 Q의 본페로니 구간(여러 번 비교할 때 넓힌 신뢰구간)이에요. 표~\ref{tab:neutral-label-contrasts}에는 청산 없이 닫은 비율에서의 대비가 있어요. 표~\ref{tab:neutral-label-recipients}에는 허락을 받는 트레이더가 있고, 표~\ref{tab:neutral-label-grid}(부록~\ref{sec:er-awp-neutral-grid})에는 트레이더 라벨 여섯 가지가 모두 있어요.

% Generated by experiments/scripts/export_latex.py -- do not edit by hand.
% Source: experiments/data/2-processed-tables-and-evaluations/revised-usd-12m/window-w1/eval_summary.json
\begin{table}[htbp]
\centering
\caption[청산 없이 닫은 비율에서 허락 쪽 점수와 송금 쪽 점수의 비교.]{허락 쪽 점수와 송금 쪽 점수를 청산 없이 닫은 비율에서 서로 비교했어요. 대상은 라벨 기간에 포지션을 세 번 이상 닫은 GMX~V2 계약 트레이더예요. 이 라벨은 두 화살표 종류 어느 것으로도 만들지 않았어요(W1의 트레이더 $n=70$, W0의 트레이더 $n=87$). 비교 쌍, 판정 대비 P, Q와 그 본페로니 97.5\% 구간, 닫은 횟수로 묶음을 나눈(층화한) 행은 이 연구가 7월부터 9월까지의 로그를 뽑기 전에 등록했어요. 이것들은 원래 단위로 한 분석을 위해 등록했고, 여기서는 고친 점수에 적용했어요. W1이 주된 기간이고, W0이 보조 기간이에요. 계약을 단위로 한 짝지은 부트스트랩이고, 다시 뽑기는 2{,}000번이에요(시드 42). 층화한 행에서는, 라벨 기간에 닫은 횟수(3--4, 5--9, 10--24, 25 이상)에 따라 묶음을 나누고 묶음마다 켄달의 $\tau_b$를 구해요. 그리고 계약 쌍의 수를 무게로 삼아 이 값들을 합쳐요. 등록한 규칙을 고친 점수에 적용한 결과는 이래요. R1(그냥 센 수의 수준에서 앞섬)은 충족되었어요. R2(페이지랭크 수준의 앞섬)는 충족되지 않았어요. R3(특이성) 규칙은 적용돼요.}
\label{tab:neutral-label-contrasts}
\fitwidth{%
\begin{tabular}{lccccc}
\toprule
대비 ($a$ 빼기 $b$) & $\tau(a)$ & $\tau(b)$ & $\Delta\tau$ & 95\% 신뢰구간 & 97.5\% 신뢰구간 \\
\midrule
\multicolumn{6}{l}{\emph{기간 W1, 2026년 7월--9월 (트레이더 $n=70$)}} \\
받은 허락 수 빼기 송금 입차수 & 0.394 & 0.253 & 0.141 & [0.037, 0.240] & [0.026, 0.256] \\
EndorseRank (활동 재시작) 빼기 AWP & 0.286 & 0.263 & 0.023 & [-0.062, 0.113] & [-0.071, 0.127] \\
C-PR ($\lambda=1$) 빼기 C-PR ($\lambda=0$) & 0.302 & 0.282 & 0.019 & [-0.093, 0.128] & -- \\
EndorseRank 빼기 AWP & 0.297 & 0.263 & 0.034 & [-0.067, 0.134] & -- \\
C-PR ($\lambda=0$) 빼기 AWP & 0.282 & 0.263 & 0.019 & [-0.031, 0.066] & -- \\
EndorseRank 빼기 받은 허락 수 & 0.297 & 0.394 & -0.097 & [-0.174, -0.009] & -- \\
\quad 받은 허락 수 빼기 송금 입차수, 닫은 횟수로 층화함 & 0.381 & 0.236 & 0.145 & [0.040, 0.328] & -- \\
\quad EndorseRank (활동 재시작) 빼기 AWP, 닫은 횟수로 층화함 & 0.273 & 0.238 & 0.034 & [-0.096, 0.171] & -- \\
\midrule
\multicolumn{6}{l}{\emph{기간 W0, 2026년 4월--6월 (트레이더 $n=87$)}} \\
받은 허락 수 빼기 송금 입차수 & 0.345 & 0.249 & 0.095 & [-0.027, 0.234] & [-0.041, 0.263] \\
EndorseRank (활동 재시작) 빼기 AWP & 0.242 & 0.254 & -0.012 & [-0.153, 0.149] & [-0.177, 0.168] \\
C-PR ($\lambda=1$) 빼기 C-PR ($\lambda=0$) & 0.237 & 0.294 & -0.057 & [-0.209, 0.092] & -- \\
EndorseRank 빼기 AWP & 0.236 & 0.254 & -0.019 & [-0.173, 0.144] & -- \\
C-PR ($\lambda=0$) 빼기 AWP & 0.294 & 0.254 & 0.040 & [-0.023, 0.098] & -- \\
EndorseRank 빼기 받은 허락 수 & 0.236 & 0.345 & -0.109 & [-0.196, -0.018] & -- \\
\quad 받은 허락 수 빼기 송금 입차수, 닫은 횟수로 층화함 & 0.342 & 0.253 & 0.089 & [-0.025, 0.252] & -- \\
\quad EndorseRank (활동 재시작) 빼기 AWP, 닫은 횟수로 층화함 & 0.229 & 0.265 & -0.036 & [-0.186, 0.145] & -- \\
\bottomrule
\end{tabular}
}
\end{table}


R1은 성립해요. W1에서 동결일의 받은 허락 수는 청산 없이 닫은 비율에서 $\tau=0.394$($[0.265,0.518]$)에 이르고, 송금 입차수는 $0.253$($[0.117,0.379]$)에 이르러요. 그래서 둘의 차이인 P는 $+0.141$(97.5\% 구간 $[0.026,0.256]$)이고, 이 구간은 0보다 위에 있어요. 닫은 횟수로 묶음을 나누면(층화하면), P는 $+0.145$($[0.040,0.328]$)예요. W0에서 P는 $+0.095$(97.5\% 구간 $[-0.041,0.263]$)로, 양수이지만 가려지지 않아요. R1은 W1의 두 구간이 모두 0보다 위에 있기를 요구했어요. 그래서 등록한 문구로 말하면, 나중에 청산을 피하는 데에서 그냥 센 수의 수준으로는 허락이 송금보다 앞서요. R2는 성립하지 않아요. Q는 AWP의 활동 재시작을 쓴 엔도스랭크에서 AWP를 뺀 값이에요(조건을 똑같이 맞춘 페이지랭크 짝). Q는 W1에서 $+0.023$, W0에서 $-0.012$이고, 두 97.5\% 구간에 모두 0이 들어 있어요($[-0.071,0.127]$와 $[-0.177,0.168]$). 그래서 페이지랭크 버전들은 어느 쪽이 나은지 가려지지 않아요. C-PR의 두 층(W1에서 $+0.019$, $[-0.093,0.128]$), 엔도스랭크 대 AWP($+0.034$, $[-0.067,0.134]$), C-PR ($\lambda=0$) 대 AWP($+0.019$, $[-0.031,0.066]$)도 어느 기간에서도 가려지지 않아요. 등록 문서가 정하지 않은 대비에서는, 엔도스랭크가 허락한 주인 수보다 낮아요. W1에서는 $0.097$($[-0.174,-0.009]$), W0에서는 $0.109$($[-0.196,-0.018]$)만큼 낮아요.

다른 라벨들은 판정에 들어가지 않아요(표~\ref{tab:neutral-label-grid}). 청산이 없었는지 나타내는 표시에서 P는 W1에서 $+0.141$($[0.044,0.232]$), W0에서 $+0.109$($[-0.022,0.257]$)예요. 활동에 따라 커지는 이익을 낸 닫기와 실현 이익에서는, W1에서 페이지랭크 짝들이 송금 쪽에 유리하게 나와요. Q($-0.215$, $[-0.399,-0.042]$와 $-0.299$, $[-0.457,-0.146]$)와 엔도스랭크 대 AWP($-0.200$, $[-0.390,-0.023]$와 $-0.292$, $[-0.460,-0.134]$)의 구간이 0보다 아래에 있어요. P도 같은 쪽을 가리키지만 가려지지는 않아요. R3은 적용돼요. W1에서 받은 허락 수는 이익을 낸 비율($-0.162$, $[-0.295,-0.025]$)과 손실이 아닌 비율($-0.222$, $[-0.328,-0.106]$)에서 송금 입차수보다 낮아요. 그래서 등록한 문구로 말하면, 이 신호는 거래를 잘하는 것이 아니라 강제 청산을 피하는 것에 관한 거예요. 손실이 아닌 비율의 차이 가운데 일부는 만든 방식 때문에 저절로 생겨요. 처리 과정은 청산된 닫기를 손실이 아닌 닫기로 세요. 그래서 청산이 있는 트레이더는 손실이 아닌 비율이 올라가요. 따라서 R3은 주로 이익을 낸 비율에 기대요. 이익을 낸 비율은 청산된 닫기를 이익을 내지 않은 것으로 세요.

% Generated by experiments/scripts/export_latex.py -- do not edit by hand.
% Source: experiments/data/2-processed-tables-and-evaluations/revised-usd-12m/window-w1/eval_summary.json
\begin{table}[htbp]
\centering
\caption[허락을 받는 트레이더와 라벨이 있는 다른 모든 트레이더의 비교.]{동결일에 받은 허락 수가 0보다 큰 GMX~V2 계약 트레이더(받는 쪽)를, 라벨이 있는 다른 모든 트레이더와 두 기간에서 비교했어요(W1의 트레이더 $n=70$, W0의 트레이더 $n=87$). 받는 쪽과 다른 트레이더의 비교, 그리고 영향 점검은 표~\ref{tab:neutral-label-contrasts}의 대비와 함께 등록했어요. 여기서는 고친 점수에 적용했어요. 트레이더 가운데 가장 큰 받은 허락 수는 W1에서 1, W0에서 3이에요. 그래서 영향 점검에서 빼는 받는 쪽 열 개는 같은 값(동점) 가운데서 점(노드)의 순서대로 골랐어요. 영향 점검은 남은 트레이더에 대해 계약을 단위로 한 짝지은 부트스트랩이에요(다시 뽑기 2{,}000번, 시드 42).}
\label{tab:neutral-label-recipients}
\fitwidth{%
\begin{tabular}{lcc}
\toprule
 & 받는 쪽 & 다른 트레이더 \\
\midrule
\multicolumn{3}{l}{\emph{기간 W1, 2026년 7월--9월 (트레이더 $n=70$)}} \\
트레이더 수 & 37 & 33 \\
청산이 없는 비율 & 1.000 & 0.727 \\
청산 없이 닫은 비율의 평균 & 1.000 & 0.878 \\
이익을 낸 비율의 평균 & 0.325 & 0.479 \\
손실이 아닌 비율의 평균 & 0.433 & 0.644 \\
라벨 기간에 닫은 횟수의 중앙값 & 24.0 & 11.0 \\
\multicolumn{3}{l}{받은 허락 수가 가장 큰 받는 쪽 열 개를 뺀 P: $\Delta\tau=0.155$ [0.032, 0.276]} \\
\midrule
\multicolumn{3}{l}{\emph{기간 W0, 2026년 4월--6월 (트레이더 $n=87$)}} \\
트레이더 수 & 39 & 48 \\
청산이 없는 비율 & 1.000 & 0.750 \\
청산 없이 닫은 비율의 평균 & 1.000 & 0.919 \\
이익을 낸 비율의 평균 & 0.359 & 0.474 \\
손실이 아닌 비율의 평균 & 0.401 & 0.607 \\
라벨 기간에 닫은 횟수의 중앙값 & 30.0 & 10.0 \\
\multicolumn{3}{l}{받은 허락 수가 가장 큰 받는 쪽 열 개를 뺀 P: $\Delta\tau=0.088$ [-0.048, 0.240]} \\
\bottomrule
\end{tabular}
}
\end{table}


그냥 센 수의 수준에서 한 비교는 몇 안 되는 트레이더에 바탕을 두어요. W1에서 트레이더들 가운데 가장 큰 받은 허락 수는 1이에요. 그래서 P는 사실상 허락을 받은 트레이더 37개와 나머지 트레이더 33개를 견주는 거예요. W0에서는 가장 큰 값이 3이고, 트레이더 87개 가운데 39개가 허락을 받은 트레이더예요. 허락을 받은 트레이더는 두 기간 모두에서 하나도 빠짐없이 청산을 피했어요(37개 중 37개, 39개 중 39개). 나머지 트레이더는 W1에서 72.7\%, W0에서 75.0\%만 청산을 피했어요. 허락을 받은 트레이더는 포지션도 더 많이 닫았어요(중앙값이 W1에서 24 대 11, W0에서 30 대 10). 그리고 이익을 낸 비율과 손실이 아닌 비율이 더 낮았어요(W1에서 손실이 아닌 비율의 평균 0.433 대 0.644). 이것이 R3 뒤에 있는 모양이에요. 받은 허락 수가 가장 큰, 허락을 받은 트레이더 10개를 빼면, P는 W1에서 $+0.155$($[0.032,0.276]$), W0에서 $+0.088$($[-0.048,0.240]$)이에요.

\dissertationSubheading[sec:empirical-claims]{주장(C1--C6)}

이 장의 결과는 아래의 주장들을 뒷받침해요. 각 주장은 표~\ref{tab:hrq-claim-map}의 가설 열과 연구 질문 열에 이어져 있어요. C1--C5가 주된 주장이에요. C2와 C3은 라벨을 본 뒤에 계산한 대비에 바탕을 두고, C6은 두 종류의 화살표 어느 쪽으로도 만들지 않은 결과 값에서의 비교에 바탕을 둬요. C1의 규칙~A와~B의 판정, 그리고 C6의 R1--R3의 판정은 등록한 규칙을 고친 점수에 재분석으로 적용한 거예요. 등록 분석도 같은 판정에 이르렀지만, R1만은 달랐어요(부록~\ref{sec:map-registered}).

\begin{enumerate}
\item[C1.] 자기 규칙 아래에서 C-PR은 한 층만 걷는 자기의 두 층을 이기지 못해요. C-PR은 달러 무게를 단 두 층을 고정한 $\lambda$로 섞고, 고르게 다시 출발하는 점수예요. 규칙~A는 실패해요. C-PR은 새 허락 쌍에서 허락 걷기보다 나쁘지 않지만($+0.031$, $[0.021,0.042]$), 새 송금자에서는 송금 걷기보다 나빠요($-0.030$, $[-0.040,-0.020]$). 그리고 같은 기간 부분도 실패해요(표~\ref{tab:holdout-diff}과~\ref{tab:tau-diff-hybrid}). 규칙~B는 충족되지 않아요(표~\ref{tab:fresh-contrasts}). 라벨을 본 뒤에 찾은, 새 허락 쌍에서 송금 걷기보다 나은 W0의 이득은 W1에서도 다시 나타나요(F1, $+0.170$, $[0.149,0.189]$). 하지만 C-PR은 새 송금자에서 송금 걷기보다 허용 폭 $0.01$을 넘게 낮아요(F2, $-0.025$, $[-0.038,-0.015]$). 그리고 F3은 아슬아슬하게 실패해요($+0.000$, $[-0.010,0.010]$). 첫 번째 민감도 분석에서 AWP와 견주면, F2는 뚜렷하게 실패해요($-0.075$, $[-0.085,-0.065]$). S-PR은 새 송금자를 송금 걷기보다 조금 더 잘 줄 세우고, 새 허락 쌍은 허락 걷기보다 훨씬 못하게 줄 세워요. 그래서 허락 화살표는 주로 걷기를 실어 나를 때 도움이 돼요. 혼합 점수들이 기간 안에서 어디에 있는지는 \ref{sec:coupled-results}절에서 보고해요.
\item[C2.] 규칙 밖의 비교에서, C-PR은 두 기간 모두에서 두 스펜더 라벨을 엔도스랭크보다 더 잘 줄 세워요(새 허락 쌍 $+0.026$, $[0.016,0.037]$와 $+0.017$, $[0.006,0.029]$; 새 송금자 $+0.064$, $[0.052,0.075]$와 $+0.064$, $[0.051,0.076]$). 그리고 기간 안에서는 송금 가족과 시빌 안정성 가족에서 엔도스랭크보다 앞서요. C-PR은 새 허락 쌍을 AWP보다 더 잘 줄 세워요($+0.133$, $[0.114,0.153]$; $+0.159$, $[0.139,0.178]$). 이것은 대부분 AWP가 0점을 주는 스펜더들 덕분이에요. 송금을 받은 적이 있는 스펜더들 사이에서는, W0에서 둘이 비슷하고 W1에서 C-PR이 $0.022$($[0.011,0.033]$)만큼 앞서요. C-PR은 두 기간 모두에서 새 송금자를 AWP와 송금 입차수보다 못하게 줄 세워요. 새 허락 쌍은 두 기간 모두에서 받은 허락 수보다 못하게 줄 세우고($-0.063$, $[-0.076,-0.052]$; $-0.065$, $[-0.077,-0.054]$), W1에서는 활동 재시작을 쓴 엔도스랭크보다도 못하게 줄 세워요($-0.019$, $[-0.029,-0.008]$). 그리고 W1 트레이더의 청산 없이 닫은 비율은 받은 허락 수보다 못하게 줄 세워요($-0.112$, $[-0.194,-0.026]$)(표~\ref{tab:cpr-contrasts}). 이 대비들 가운데 어느 것도 미리 정해 두지 않았어요.
\item[C3.] 기간 밖에서는 어느 한 층 걷기도 자기가 매끄럽게 다듬는 원시 차수를 이기지 못해요. AWP는 두 기간 모두에서 나중의 새 송금자를 송금 입차수보다 못하게 줄 세워요(W0 $-0.038$, $[-0.043,-0.032]$; W1 $-0.041$, $[-0.047,-0.034]$). 엔도스랭크는 두 기간 모두에서 나중의 새 허락 쌍을 받은 허락 수보다 못하게 줄 세워요(W0 $-0.090$, $[-0.100,-0.079]$; W1 $-0.082$, $[-0.090,-0.073]$)(표~\ref{tab:holdout-diff}과~\ref{tab:fresh-posthoc}). W0에서는 두 비교 모두 점수가 닿는 범위의 문제가 아니에요(표~\ref{tab:holdout-receiving}). 그리고 W1에서는 이 두 대비를 이렇게 좁혀서 계산하지 않았어요. 활동 재시작을 써도 엔도스랭크는 그냥 센 수보다 낮아요(\ref{sub:endorserank-restarts}절). 그래서 이 라벨들에서는 다시 출발 규칙이 무엇이든, 퍼뜨림이 같은 종류 화살표를 센 수에 아무것도 보태지 않아요. W0 계약 트레이더의 이익 라벨에서 짝지은 대비가 두 쪽을 갈라 놓는 것은 실현 이익뿐이에요. 실현 이익에서는 송금 층만 걷는 것이 허락 층보다 더 잘 줄 세워요(C-PR ($\lambda=1$)에서 C-PR ($\lambda=0$)을 뺀 값 $-0.257$, $[-0.447,-0.071]$). 이익을 낸 닫기에서는 갈라 놓지 못해요(엔도스랭크에서 AWP를 뺀 값 $-0.046$, $[-0.205,0.113]$). 그리고 어떤 점수도 닫기 가운데 이익을 낸 비율을 줄 세우지 못해요(표~\ref{tab:holdout-traders}와~\ref{tab:neutral-label-grid}). 이 비교들은 라벨을 본 뒤에 계산했어요.
\item[C4.] 같은 기간에서 각 한 층 걷기는 자기가 걷는 화살표 종류의 대리 지표와 가장 잘 맞아요. 엔도스랭크는 허락 가족과($0.560$, $[0.548,0.571]$), AWP는 송금 가족과($0.736$, $[0.728,0.744]$) 가장 잘 맞아요. 그리고 C-PR의 섞은 점수들은 자기의 두 층 사이에 있어요(표~\ref{tab:alignment-hybrid}). 이 시험들은 내적 일관성(같은 자료 안에서 앞뒤가 맞는지)을 재요. 그리고 그 일치의 일부는 만든 방식 때문에 저절로 생겨요. 외적 타당도(바깥 자료로도 맞는지)는 이 시험들이 다루는 범위 밖이에요. 두 걷기는 코호트를 서로 다르게 줄 세워요. 둘의 일치 $\tau_b=0.199$($[0.181,0.221]$)의 구간에는 0도 1도 들어 있지 않아요. 계약 쌍의 $96.4\%$는 두 점수 모두에서 순서가 정해져 있고, 그 가운데 $39.8\%$는 순서가 어긋나는 쌍이에요. 어느 결과도 동점이나, 그래프 밖에 있는 계약을 다루는 규칙에서 나오지 않아요(표~\ref{tab:tie-sensitivity}, 그림~\ref{fig:rank-scatter}).
\item[C5.] 똑같은 계약 집합에서, 각 점수의 풀이 프로그램을 한 번 부르는 비용은 그 그래프의 크기와 함께 커지고, 그 걷기에 필요한 반복 횟수와 함께 커져요. 반복 한 번에 드는 시간이 $|E|$와 $|V|$에 따라 정해지기 때문이에요. 허락 그래프의 화살표 수는 송금 그래프의 45.2\%예요. 그리고 AWP는 엔도스랭크보다 $3.21$배 오래 걸리고, 이때 화살표 수 비율은 $2.21$이에요(표~\ref{tab:benchmark-runtime}). 부분 표본에서는 어느 단계에서도 시간 비율이 화살표 수 비율을 따르지 않아요. 그리고 단계들은 코호트 그래프보다 화살표가 많은 경우를 시험하지 않아요(표~\ref{tab:benchmark-scaling}). 비용이 낮은 것은 알고리즘의 장점이 아니에요. C-PR과 S-PR은 송금 층을 걷고, AWP보다 오래 걸려요. C-PR은 $1.36$배 오래 걸려요.
\item[C6.] GMX~V2 계약 트레이더가 나중에 청산을 피하는 일은, 두 종류의 화살표 어느 쪽으로도 만들지 않은 결과 값 가운데 거래 상대의 위험(거래 상대 때문에 손실을 볼 위험)과 관계있는 단 하나예요. 이 결과 값에서 동결일의 허락한 주인 수는 주된 기간 W1에서 송금 수보다 앞서요(P $=+0.141$, 97.5\% 구간 $[0.026,0.256]$). 반면에 W0에서 P는 양수이지만 가려지지 않아요($+0.095$, 97.5\% 구간 $[-0.041,0.263]$). 페이지랭크 버전들은 가려지지 않아요(Q $=+0.023$, 97.5\% 구간 $[-0.071,0.127]$). W1에서 허락한 주인 수는, 닫기 가운데 이익을 낸 비율과 손실이 아닌 비율에서 트레이더를 송금 수보다 낮게 줄 세워요. 그래서 R3의 문구로 말하면, 이 신호는 거래를 잘하는 것이 아니라 강제 청산을 피하는 것에 관한 거예요. W0에서는 어느 차이도 가려지지 않아요(표~\ref{tab:neutral-label-grid}). 표본이 작아요. 계약 트레이더는 70개와 87개이고, 모든 구간이 넓어요(\ref{sub:neutral-label-results}절).
\end{enumerate}

같은 기간 표들과 순위 갈라짐은 RQ1에 답해요. 벤치마크 표들은 RQ2에, 한 층 걷기의 홀드아웃들은 표~\ref{tab:cpr-contrasts}과 함께 RQ3에, \ref{sec:cpr-results}절은 RQ4와 RQ5에 답해요. C6은 연구 질문을 더하지 않아요. 엔도스랭크와 AWP를 직접 견준 비교는 부록~\ref{ch:appendix-er-awp}와 표~\ref{tab:neutral-label-contrasts}에 있어요. 그 가운데 어느 것도 이 논문의 주장을 담고 있지 않아요. \ref{ch:discussion}장에서는 이 발견들을 \ref{ch:introduction}장의 연구 질문과 목표에 비추어 읽어요.
